% !TEX TS-program = xelatex
% Command for running this example (needs latexmkrc file):
%    latexmk -bibtex -pdf main.tex

%	نمونه پایان‌نامه آماده شده با استفاده از کلاس tehran-thesis، نگارش 1
%	سینا ممکن، دانشگاه تهران 
%	https://github.com/sinamomken/tehran-thesis
%	گروه پارسی‌لاتک
%	http://www.parsilatex.com
%	این نسخه، بر اساس نسخه‌ 0.1 از کلاس IUST-Thesis آقای محمود امین‌طوسی آماده شده است.
%        http://profsite.sttu.ac.ir/mamintoosi

%----------------------------------------------------------------------------------------------
% اگر قصد نوشتن پروژه کارشناسی را دارید، در خط زیر به جای msc، کلمه bsc و اگر قصد نوشتن رساله دکترا را دارید، کلمه phd را قرار دهید. کلیه تنظیمات لازم، به طور خودکار، اعمال می‌شود.

% اگر مایلید پایان‌نامه شما دورو باشد به جای oneside در خط زیر از twoside استفاده کنید.

% برای حاشیه‌نویسی و کم کردن صفحات ابتدایی، گزینه draft را وارد و برای نسخه نهایی آن را حذف کنید.

% برای استفاده از قلم‌های سری IR Series گزینه irfonts را وارد و برای استفاده از قلم‌های X Series 2 آن را حذف کنید.

\documentclass[
twoside
% ,openany
,bsc
,irfonts
% ,draft
]{./tex/tehran-thesis}

% فایل commands.tex را مطالعه کنید؛ چون دستورات مربوط به فراخوانی بسته‌ها، فونت و دستورات خاص در این فایل قرار دارد.
\input{./tex/commands}

% مشخصات پایان‌نامه را در فایلهای faTitle و enTitle وارد نمایید.
% !TeX root=../main.tex
% در این فایل، عنوان پایان‌نامه، مشخصات خود، متن تقدیمی‌، ستایش، سپاس‌گزاری و چکیده پایان‌نامه را به فارسی، وارد کنید.
% توجه داشته باشید که جدول حاوی مشخصات پروژه/پایان‌نامه/رساله و همچنین، مشخصات داخل آن، به طور خودکار، درج می‌شود.
%%%%%%%%%%%%%%%%%%%%%%%%%%%%%%%%%%%%
% دانشگاه خود را وارد کنید
\university{دانشگاه تهران}
% پردیس دانشگاهی خود را اگر نیاز است وارد کنید (مثال: فنی، علوم پایه، علوم انسانی و ...)
\college{پردیس دانشکده‌های فنی}
% دانشکده، آموزشکده و یا پژوهشکده  خود را وارد کنید
% TODO تأیید کنید — از قالب رسمی BScThesis_template_new.docx همین پروژه استخراج شده
\faculty{دانشکده مهندسی برق و کامپیوتر}
% گروه آموزشی خود را وارد کنید (در صورت نیاز)
% TODO — نام دقیق گروه را وارد کنید (این مقدار placeholder است تا سند کامپایل شود)
\department{TODO}
% رشته تحصیلی خود را وارد کنید
% TODO تأیید کنید
\subject{مهندسی کامپیوتر}
% گرایش خود را وارد کنید
% TODO — نام دقیق گرایش را وارد کنید (مثلاً هوش مصنوعی؛ placeholder تا تکمیل)
\field{TODO}
% عنوان پایان‌نامه را وارد کنید
% TODO تأیید یا اصلاح کنید — پیشنهاد بر اساس AGENTS.md (تعریف پروژه)
\title{پیش‌بینی نرخ عدم‌دریافت غذا در سلف‌های دانشگاه تهران با رگرسیون کوانتایل روی
سری‌های زمانی چندگانه}
% نام استاد(ان) راهنما را وارد کنید
% TODO — نام و مرتبه‌ی استاد راهنما را وارد کنید
\firstsupervisor{}
\firstsupervisorrank{}
% اگر استاد راهنمای دوم ندارید، دو خط زیر را غیرفعال کنید.
%\secondsupervisor{دکتر راهنمای دوم}
%\secondsupervisorrank{استادیار}
% نام استاد(دان) مشاور را وارد کنید. چنانچه استاد مشاور ندارید، دستورات پایین را غیرفعال کنید.
%\firstadvisor{دکتر مشاور اول}
%\firstadvisorrank{استادیار}
%\secondadvisor{دکتر مشاور دوم}
% نام داوران داخلی و خارجی — نزدیک دفاع تکمیل می‌شود، فعلاً لازم نیست.
%\internaljudge{دکتر داور داخلی}
%\internaljudgerank{دانشیار}
%\externaljudge{دکتر داور خارجی}
%\externaljudgerank{دانشیار}
%\externaljudgeuniversity{دانشگاه داور خارجی}
% نام نماینده کمیته تحصیلات تکمیلی در دانشکده \ گروه — نزدیک دفاع تکمیل می‌شود.
%\graduatedeputy{دکتر نماینده}
%\graduatedeputyrank{دانشیار}
% نام دانشجو را وارد کنید
% TODO
\name{}
% نام خانوادگی دانشجو را وارد کنید
% TODO
\surname{}
% شماره دانشجویی دانشجو را وارد کنید
% TODO
\studentID{}
% تاریخ پایان‌نامه را وارد کنید
% TODO
\thesisdate{}
% به صورت پیش‌فرض برای پایان‌نامه‌های کارشناسی تا دکترا به ترتیب از عبارات «پروژه»، «پایان‌نامه» و «رساله» استفاده می‌شود؛ اگر  نمی‌پسندید هر عنوانی را که مایلید در دستور زیر قرار داده و آنرا از حالت توضیح خارج کنید.
%\projectLabel{پایان‌نامه}

% به صورت پیش‌فرض برای عناوین مقاطع تحصیلی کارشناسی تا دکترا به ترتیب از عبارت «کارشناسی»، «کارشناسی ارشد» و «دکتری» استفاده می‌شود؛ اگر نمی‌پسندید هر عنوانی را که مایلید در دستور زیر قرار داده و آنرا از حالت توضیح خارج کنید.
%\degree{}
%%%%%%%%%%%%%%%%%%%%%%%%%%%%%%%%%%%%%%%%%%%%%%%%%%%%
%% پایان‌نامه خود را تقدیم کنید! (اختیاری) %%
%% متن نمونه‌ی قالب اصلی (خانواده‌ی نگارنده‌ی قالب) غیرفعال شد — این صفحه اختیاری است؛
%% اگر می‌خواهید، متن خودتان را اینجا بنویسید و دستور \dedication را فعال کنید.
%\dedication
%{
%{\Large تقدیم به:}\\
%\begin{flushleft}{
%	\huge
%	...
%}
%\end{flushleft}
%}
%% متن قدردانی (اختیاری) %%
%% متن نمونه‌ی قالب اصلی غیرفعال شد — اگر می‌خواهید، متن خودتان (شامل نام استاد راهنما)
%% را اینجا بنویسید و دستور \acknowledgement را فعال کنید.
%\acknowledgement{
%...
%}
%%%%%%%%%%%%%%%%%%%%%%%%%%%%%%%%%%%%
%چکیده پایان‌نامه را وارد کنید
% TODO — چکیده‌ی فارسی (حداکثر ۱ صفحه) در فاز نگارش (پس از تکمیل فصل‌های ۲ تا ۵) نوشته
% می‌شود؛ نه پیش از آن. منبع محتوا: doc/project-definition.md + نتایج نهایی فصل ۴.
\fa-abstract{
}
% کلمات کلیدی پایان‌نامه را وارد کنید
% TODO تأیید یا اصلاح کنید
\keywords{پیش‌بینی تقاضا، رگرسیون کوانتایل، مسئله‌ی روزنامه‌فروش، سری‌های زمانی، هدررفت غذا}
% انتهای وارد کردن فیلد‌ها
%%%%%%%%%%%%%%%%%%%%%%%%%%%%%%%%%%%%%%%%%%%%%%%%%%%%%%

% مشخصات انگلیسی پایان‌نامه
% !TeX root=../main.tex
% در این فایل، عنوان پایان‌نامه، مشخصات خود و چکیده پایان‌نامه را به انگلیسی، وارد کنید.

%%%%%%%%%%%%%%%%%%%%%%%%%%%%%%%%%%%%
\latinuniversity{University of Tehran}
\latincollege{College of Engineering}
% TODO تأیید کنید — معادل انگلیسی «دانشکده مهندسی برق و کامپیوتر»
\latinfaculty{School of Electrical and Computer Engineering}
% عمداً خالی گذاشته شده
\latindepartment{}
\latinsubject{Computer Engineering}
\latinfield{Hardware}
% TODO تأیید کنید — ترجمه‌ی انگلیسی، معادل رسمی داده نشده بود
\latintitle{Optimizing Food Preparation Quantity Prediction for University of Tehran
Dining Halls Using Time Series Models}
\firstlatinsupervisor{Babak Najar Araabi}
%\secondlatinsupervisor{Second Supervisor}
%\firstlatinadvisor{First Advisor}
%\secondlatinadvisor{Second Advisor}
\latinname{Mahdi}
\latinsurname{Vajhi}
\latinthesisdate{September 2026}
\latinkeywords{Food Demand Forecasting, Machine Learning, No-show Rate, Asymmetric
Loss, Time Series, Food Waste Reduction}
% ترجمه‌ی چکیده‌ی فارسی در tex/faTitle.tex — اعداد هر دو باید هم‌زمان به‌روز شوند.
\en-abstract{
At the University of Tehran dining halls, the amount of food cooked for each meal is
set from the number of registered reservations, yet catering-system data covering Azar
1402 to Khordad 1403 (November 2023 to June 2024) shows that 8.05\% of reservations are
never collected: more than 162,000 portions in five months, close to 61 billion tomans
at today's prices (September 2026; prices were lower during the data period).
Cooking less reduces this waste but raises the risk of a shortage, and the two errors do
not carry the same cost.

This study estimates the no-show rate ahead of each meal and converts it into a cooking
quantity. Instead of estimating the conditional mean, the problem is formulated as
estimating the $\tau$-th quantile of the conditional distribution of that rate at the
(day, meal, dining hall, food type) level, so that the asymmetry between shortage and
surplus costs enters the objective function itself. Predictions are issued at 15:00 for the next day's lunch and at
23:00 for the next day's dinner, and no feature that is not yet available at that moment
enters the model. The raw data consist of 2,564,919 individual reservations, which
after cleaning were aggregated into 7,579 cells over 142 serving days and 30 dining
halls; one hundred features were built on this basis and thirteen model families were
screened.

The twenty-five surviving models were evaluated with the Pinball loss at $\tau = 0.20$
on five expanding time folds and one locked holdout window, against eight baselines.
Only five models significantly beat the strongest baseline, a simple lookup table over
dining hall and day of week, and the best of them reduced the loss by about nine
percent. On the locked window, which spans Ramadan and the days of mourning after
the presidential helicopter crash, the ranking changed and only the Bayesian
hierarchical model retained a significant edge.

The quantile level $\tau$ is left to dining hall management as a policy lever that sets
the accepted level of risk. At the operating point
$\tau = 0.20$, waste falls by about 71\% relative to cooking the full reservation
count, at the cost of an average shortfall of 31 portions per day across the whole
university, which is 0.28\% of total demand; the three policy scenarios examined yield,
at today's prices, annual savings between 36.3 and 61.9 billion tomans. These figures are
conditional on explicit cost assumptions, have not been tested in a live setting, and do
not generalize to the summer or examination periods.
}



% تنظیمات و تعاریف واژه‌نامه و اختصارات
\input{./tex/glossaries-settings}

% واژه‌نامه‌ی فارسی-انگلیسی این گزارش.
% نمونه‌های قالب اصلی (Gloss, Acronym, Description, Draft, Absorption, RandomVariable,
% Action, Optimization) که با موضوع پروژه بی‌ربط بودند حذف شدند.
%
% TODO: اصطلاحات لاتینی که هنگام نگارش فصل‌های ۲ و ۳ اولین‌بار ترجمه می‌شوند، اینجا با
% دستور \newword{برچسب}{English}{ترجمه‌ی فارسی مفرد}{ترجمه‌ی فارسی جمع} اضافه شوند.
% فهرست نامزد (doc/report-outline.md بند ۳): Newsvendor، Pinball Loss،
% Quantile Regression، Conformal Prediction، Overdispersion، Rolling-origin،
% Data Leakage، Calibration، Coverage، Surrogate Model، Feature Importance.
%
% نمونه:
% \newword{Newsvendor}{Newsvendor}{روزنامه‌فروش}{}

% فهرست علائم اختصاری این گزارش.
% نمونه‌های فیزیکی قالب اصلی (a=شتاب، F=نیرو) که با موضوع پروژه بی‌ربط بودند حذف شدند.
%
% TODO: اختصارات این پروژه (doc/report-outline.md بند ۳) اینجا با دستور
% \newacronym{برچسب}{مخفف}{شرح فارسی} اضافه شوند:
% CQR, CV, DVC, GBM, MAE, RMSE, SHAP, PDP, ALE, ADF, KPSS, ARCH-LM, DM, FDR,
% AQI, PSI, MLOps, WBS.
%
% نمونه:
% \newacronym{CQR}{\lr{CQR}}{\rl{رگرسیون کوانتایل همنوا}}


\begin{document}

\pagenumbering{adadi} % یک، دو، ...
% ابتدای درج صفحات مختلف
\coverPage
% بررسی حالت پیش‌نویس
\ifoptiondraft{}{%
    \besmPage
    \titlePage
%%%%%%%%%%%%%%%%%%%%%%%%%%%
    \esalatPage
% چنانچه مایل به چاپ صفحات «تقدیم»، «نیایش» و «سپاس‌گزاری» در خروجی نیستید، خط‌های زیر را با گذاشتن ٪  در ابتدای آنها غیرفعال کنید.
    \taghdimPage
    \ghadrdaniPage
} % end ifoptiondraft
\abstractPage
% شروع درج فهرست‌ها
\newpage\cleardoublepage
\pagenumbering{harfi} % آ، ب، ...
% فهرست‌های زیر عمداً تک‌فاصله‌اند — \doublespacing سراسری commands.tex (برای بدنه‌ی
% متن) بدون این override روی ردیف‌های فهرست مطالب/تصاویر/جداول هم اعمال می‌شد و فاصله‌ی
% هر ردیف را به‌اندازه‌ی یک پاراگراف بدنه می‌کرد.
\begin{singlespace}
\tableofcontents \clearpage
% بررسی حالت پیش‌نویس برای بقیه فهرست‌ها
\ifoptiondraft{
    \listoftodos
}{%
    \listoffigures \clearpage
    \listoftables  \clearpage
    \printacronyms
} % end ifoptiondraft
\end{singlespace}

\pagestyle{fancy}
\pagenumbering{arabic} % 1, 2, ...

% !TeX root=../main.tex

\chapter{مقدمه و بیان مسئله}

\section{مقدمه}
\label{sec:1-1}

در سلف‌های دانشگاه تهران غذا بر مبنای تعداد رزروهای ثبت‌شده پخته می‌شود، ولی بخشی از این
رزروها هرگز تحویل گرفته نمی‌شوند و غذای متناظرشان، که هزینه‌اش پرداخت شده، مصرف‌کننده‌ای
ندارد. کاهش این هدررفت با پخت کمتر، بدون تخمین اینکه چقدر کمتر، ریسک کمبود غذا را می‌سازد
که هزینه‌اش برای دانشگاه سنگین‌تر است. مسئله بنابراین یک مسئله‌ی پیش‌بینی است و نه صرفاً یک
مسئله‌ی عملیاتی: پیش از هر وعده باید تخمین زد چه نسبتی از رزروها دریافت نخواهد شد. این
پژوهش همان نسبت را با رگرسیون کوانتایل روی سری‌های زمانی چندگانه برآورد می‌کند و آن را به
عددی برای پخت تبدیل می‌کند که هزینه‌ی نابرابر کمبود و مازاد در آن لحاظ شده است.

\section{تاریخچه‌ای از موضوع تحقیق}
\label{sec:1-2}

تخصیص غذا در دانشگاه تهران از طریق یک سامانه‌ی رزرو اینترنتی انجام می‌شود. دانشجو غذای هر
وعده را دست‌کم ۷۲ ساعت پیش از سرو رزرو می‌کند و پس از بسته‌شدن مهلت امکان لغو ندارد؛ پس در
لحظه‌ای که آشپزخانه باید درباره‌ی مقدار پخت تصمیم بگیرد، تعداد رزرو عددی قطعی و کاملاً معلوم
است. هر رزرو به ترکیبی از روز، وعده، سلف و نوع غذا نسبت داده می‌شود و هنگام سرو، تحویل با
کارت دانشجویی یا با کد ثبت می‌شود. رویه‌ی کنونی مدیریت پخت همین عدد رزرو را مبنا می‌گیرد؛
مصاحبه‌های آغاز پروژه نشان داد پیمانکاران در عمل گاهی بر پایه‌ی تجربه‌ی شخصی کمتر از آن
می‌پزند، بی‌آنکه این کاهش با کارفرما هماهنگ شود یا در صورتحساب بازتاب یابد. تصمیم واقعی
پخت بنابراین بیرون از سازوکار رسمی و بدون مبنای مستند گرفته می‌شود.

\section{شرح مسئله تحقیق}
\label{sec:1-3}

داده‌ی سامانه‌ی تغذیه در بازه‌ی آذر ۱۴۰۲ تا خرداد ۱۴۰۳، شامل ۱۴۲ روز سرو در ۳۰ سلف و شش
شهر، نشان می‌دهد $8.05\%$ کل رزروهای ثبت‌شده تحویل گرفته نشده‌اند. این عدد نسبت مجموع
رزروهای دریافت‌نشده به مجموع کل رزروهاست؛ میانگین ساده‌ی نرخ سلول‌ها $9.59\%$ است، یعنی
حدود ۱۹ درصد بالاتر، چون به سلف‌های کوچک و پرنوسان وزنی بیش از سهمشان در حجم غذا می‌دهد.
حجم مطلق هدررفت بیش از ۱۶۲ هزار پرس در پنج ماه است که با قیمت تقریبی ۱۲۰ هزار تومان برای
هر پرس، نزدیک به ۱۹ میلیارد تومان در همین بازه و سالانه‌شده حدود ۵۰ میلیارد تومان می‌شود.

هزینه‌ی خطا میان سه گروه درگیر یکسان توزیع نمی‌شود: مدیریت سلف مقدار پخت را تعیین می‌کند،
پیمانکار خرید و پخت را بر همان عدد می‌بندد، و دانشجو که در تصمیم نقشی ندارد پیامد خطای رو
به پایین را به شکل کمبود غذا تجربه می‌کند. پخت بیش از نیاز هزینه‌ی مالی دارد و قابل جبران
است، ولی پخت کمتر از نیاز به قطع خدمت می‌انجامد. همین عدم‌تقارن، نه دقت متقارن پیش‌بینی،
شکل تابع هدف این پژوهش را تعیین می‌کند.

\section{تعریف موضوع تحقیق}
\label{sec:1-4}

واحد تحلیل چهارتایی $(d,m,r,f)$ است: روز، وعده، سلف و نوع غذا، یعنی همان سطحی که آشپزخانه
برای آن یک عدد پخت می‌گیرد. کمیت پیش‌بینی‌شونده نرخ عدم‌دریافت $\rho$ است، نسبت رزروهای
دریافت‌نشده به کل رزروهای همان سلول، و آنچه برآورد می‌شود میانگین شرطی این نرخ نیست بلکه
کوانتایل $\tau$اُم توزیع شرطی آن است؛ دلیل هر دو انتخاب در بندهای \ref{sec:2-2-3} و
\ref{sec:3-2-1} آمده است. همین کمیت را می‌توان در سطوح مختلفی از تجمیع مدل کرد، از سطح
سلول تا سطح تک‌رزرو هر دانشجو، و این پژوهش سطح داده را خودش یک محور آزمایش قرار داد تا اثر
الگوریتم و اثر سطح تجمیع در مقایسه‌ها آمیخته نشوند. قید حاکم بر کل کار، قاعده‌ی برش
اطلاعاتی است: هر پیش‌بینی تنها با اطلاعاتی ساخته می‌شود که در لحظه‌ی واقعی تصمیم در دسترس
بوده‌اند. این لحظه برای ناهار و شام یکی نیست و تفاوتشان پیامدی دارد که در نگاه اول دیده
نمی‌شود، چنان‌که بند \ref{sec:3-2-2} نشان می‌دهد.

\section{اهداف و آرمان‌های کلی تحقیق}
\label{sec:1-5}

هدف پژوهش برآورد مقدار قابل‌پخت هر وعده در سطح $(d,m,r,f)$ است، به‌گونه‌ای که هدررفت مالی
نسبت به رویه‌ی کنونی کاهش یابد بدون آنکه ریسک کمبود غذا بالا برود، و سنجش آن با زیان
\lr{Pinball} در برابر خط پایه‌ای انجام می‌شود که خودِ پروژه می‌سازد، چون هیچ خط پایه‌ی
معتبری از پیش موجود نبود. انتخاب این معیار به‌جای \lr{RMSE} از همان عدم‌تقارن هزینه می‌آید:
\lr{RMSE} با میانگین شرطی کمینه می‌شود و جهت خطا برایش تفاوتی ندارد، در حالی که اینجا
کم‌برآوردکردن نرخ عدم‌دریافت به هدررفت و بیش‌برآوردکردنش به کمبود غذا می‌انجامد. \lr{RMSE}،
\lr{MAE} و ضریب تعیین نگه داشته شده‌اند، اما تنها برای مقایسه با ادبیات موجود و نه به‌عنوان
مبنای انتخاب مدل.

محدوده‌ی کار با چند فرض صریح بسته شده است: منوی هر وعده در لحظه‌ی رزرو معلوم است، تعداد رزرو
پس از بسته‌شدن مهلت تغییر نمی‌کند، پخت تک‌مرحله‌ای است و امکان پخت تکمیلی پس از دیدن تقاضا
وجود ندارد، ظرفیت آشپزخانه قید فعالی نیست، و سرنوشت غذای مازاد و سیاست جریمه‌ی عدم‌دریافت
بیرون از محدوده‌ی پژوهش‌اند. اثرگذارترین محدودیت اما پوشش زمانی داده است: پنج ماه آذر تا
خرداد برای الگوی هفتگی کافی است ولی هیچ اطلاعاتی درباره‌ی تابستان، شروع ترم و دوره‌ی
امتحانات ندارد، پس در این گزارش هیچ ادعای تعمیمی به آن بازه‌ها مطرح نمی‌شود.

\section{روش انجام تحقیق}
\label{sec:1-6}

کار در یازده فاز اجرا شد که شکل \ref{fig:phase-flow} ترتیب و وابستگی آن‌ها را نشان می‌دهد.
مسیر اصلی از پاکسازی و یکپارچه‌سازی داده‌ی خام آغاز می‌شود، از کاوش داده و مهندسی ویژگی
می‌گذرد، پس از تثبیت پروتکل اعتبارسنجی زمانی وارد مدل‌سازی می‌شود، و با ارزیابی،
تفسیرپذیری و ساخت لایه‌ی تصمیم پایان می‌یابد. این توالی خطی نیست: یافته‌های ارزیابی
می‌توانستند فهرست ویژگی‌ها را تغییر دهند و بازگشت به فازهای پیشین بخشی از طراحی بود، به
شرطی که مجموعه‌ی آزمون قفل‌شده در این رفت‌وبرگشت‌ها دست‌نخورده بماند.

میان فازها پنج دروازه‌ی کیفیت قرار گرفت که جدول \ref{tab:gates} اثر عملی هر کدام را
می‌آورد. هر دروازه یک ریسک روش‌شناختی مشخص را پوشش می‌داد: پیش‌رفتن با داده‌ای که هنوز
آزمون‌های سازگاری‌اش را پاس نکرده، ساختن ویژگی پیش از شناختن ساختار داده، آموزش مدل پیچیده
بدون مرجعی که بتوان برتری را نسبت به آن سنجید، و باز کردن زودهنگام مجموعه‌ی آزمون.

\begin{figure}[!ht]
\centering
\resizebox{\textwidth}{!}{%
\begin{tikzpicture}[
  font=\small,
  ph/.style={draw, rounded corners=2pt, minimum width=2.6cm, minimum height=1.05cm,
             align=center, fill=black!4},
  crit/.style={ph, fill=black!14, thick},
  gate/.style={draw, rounded corners=1pt, fill=black!70, text=white,
               font=\footnotesize\bfseries, inner sep=2pt},
  ar/.style={->, thick},
  fb/.style={->, dashed, gray}
]
% ردیف اول: راست به چپ
\node[ph]   (p0) at (11.4,5) {\rl{فاز ۰}\\\rl{زیرساخت}};
\node[ph]   (p1) at (7.6,5)  {\rl{فاز ۱}\\\rl{تعریف مسئله}};
\node[ph]   (p2) at (3.8,5)  {\rl{فاز ۲}\\\rl{جمع‌آوری داده}};
\node[crit] (p3) at (0,5)    {\rl{فاز ۳}\\\rl{پاکسازی}};
% ردیف دوم: چپ به راست
\node[crit] (p4) at (0,3)    {\rl{فاز ۴}\\\rl{کاوش داده}};
\node[crit] (p5) at (3.8,3)  {\rl{فاز ۵}\\\rl{مهندسی ویژگی}};
\node[ph]   (p6) at (7.6,3)  {\rl{فاز ۶}\\\rl{پروتکل ارزیابی}};
% ردیف سوم: راست به چپ
\node[crit] (p7) at (11.4,1) {\rl{فاز ۷}\\\rl{مدل‌سازی}};
\node[crit] (p8) at (7.6,1)  {\rl{فاز ۸}\\\rl{ارزیابی}};
\node[ph]   (p9) at (3.8,1)  {\rl{فاز ۹}\\\rl{تفسیرپذیری}};
\node[ph]   (p10) at (0,1)   {\rl{فاز ۱۰}\\\rl{لایه‌ی تصمیم}};
% ردیف چهارم
\node[ph, minimum width=8cm] (p11) at (5.7,-1.4) {\rl{فاز ۱۱ — مستندسازی و ارائه (موازی با همه‌ی فازها)}};

\draw[ar] (p0) -- (p1);
\draw[ar] (p1) -- (p2);
\draw[ar] (p2) -- (p3);
\draw[ar] (p3) -- (p4) node[gate, midway] {\lr{M1}};
\draw[ar] (p4) -- (p5) node[gate, midway] {\lr{M2}};
\draw[ar] (p5) -- (p6);
\draw[ar] (p6) -| (p7) node[gate, pos=0.42] {\lr{M3}};
\draw[ar] (p7) -- (p8) node[gate, midway] {\lr{M4}};
\draw[ar] (p8) -- (p9);
\draw[ar] (p9) -- (p10);
\draw[ar] (p10) |- (p11) node[gate, pos=0.30] {\lr{M5}};
\draw[fb] (p8) -- (p5);
\end{tikzpicture}}
\caption{جریان یازده فاز پروژه و پنج دروازه‌ی کیفیت. کادرهای پررنگ‌تر مسیر بحرانی‌اند؛
پیکان خط‌چین نمونه‌ای از مسیرهای بازگشت مجاز (ارزیابی $\rightarrow$ مهندسی ویژگی) است
که ماهیت حلقوی فرآیند را نشان می‌دهد.}
\label{fig:phase-flow}
\end{figure}

\begin{table}[!ht]
\centering
\singlespacing
\caption{دروازه‌های کیفیت و اثر عملی هر کدام}
\label{tab:gates}
\vspace{2mm}
\begin{tabular}{|p{1.3cm}|p{2.4cm}|p{9cm}|}
\hline
\textbf{دروازه} & \textbf{پایان فاز} & \textbf{چه چیزی را متوقف کرد} \\
\hline
\lr{M1} & ۳ (پاکسازی) & تا وقتی یک اسکریپت واحد داده‌ی خام را به مجموعه‌ی پردازش‌شده نرساند و آزمون‌های تقاطعی و هش \lr{snapshot} ثبت نشوند، هیچ ویژگی‌ای ساخته نشد \\
\hline
\lr{M2} & ۴ (کاوش داده) & پیش از پرشدن دفتر حقایق داده و تصمیم درباره‌ی سراسری یا خوشه‌ای بودن مدل، فاز ویژگی شروع نشد \\
\hline
\lr{M3} & ۶ (پروتکل ارزیابی) & مؤثرترین دروازه: تا پیش از اجرای همه‌ی خط پایه‌ها و تثبیت پروتکل اعتبارسنجی زمانی، هیچ مدل پیچیده‌ای آموزش ندید؛ بدون مرجع مقایسه، برتری هر مدل تنها با معیاری سرخود قابل‌ادعا بود \\
\hline
\lr{M4} & ۷ (مدل‌سازی) & هر پیکربندی وضعیت مستندی گرفت («تکمیل‌شده» یا «ناسازگار-مستند») و هر اجرا با هش \lr{fold} یا \lr{snapshot} نامنطبق از مقایسه کنار رفت؛ پنجره‌ی انتهایی در کل فاز دست‌نخورده ماند \\
\hline
\lr{M5} & ۸ تا ۱۰ و گزارش & مجموعه‌ی آزمون فقط یک بار باز شد؛ پس از آن هیچ تغییری در مدل انجام نشد \\
\hline
\end{tabular}
\end{table}

\section{ساختار پایان‌نامه}
\label{sec:1-7}

فصل دوم مفاهیم لازم برای خواندن بقیه‌ی گزارش را می‌آورد: مسئله‌ی روزنامه‌فروش و زیان
نامتقارن، رگرسیون کوانتایل و پیش‌بینی احتمالاتی، اعتبارسنجی و نشت در داده‌ی زمانی،
خانواده‌های مدل، ابزارهای آماری به‌کاررفته، و مروری بر پیشینه. فصل سوم روش پیشنهادی را
تشریح می‌کند، از صورت‌بندی مسئله و قاعده‌ی برش تا مهندسی ویژگی، غربالگری مدل‌ها،
کالیبراسیون و لایه‌ی تصمیم، و در کنار آن منابع داده، پاکسازی، یافته‌های کاوش و ممیزی نشت
را گزارش می‌کند. فصل چهارم همان روش را می‌سنجد: محیط اجرا، معیارها و پروتکل اعتبارسنجی،
خط پایه‌ها، نتایج روی مجموعه‌ی قفل‌شده، و تحلیل خطا، تفسیرپذیری و اثر کسب‌وکاری. فصل پنجم
جمع‌بندی، محدودیت‌ها و پیشنهادهای کار آینده را می‌آورد و فصل ششم فهرست مراجع است.
		% فصل ۱: مقدمه و بیان مسئله
% !TeX root=../main.tex
\chapter{مروری بر مطالعات انجام شده}

% TODO ۲-۱ مقدمه فصل — یک پاراگراف کوتاه درباره‌ی چارچوب مرور: نظری (Newsvendor) + صنایع
% آنالوگ + کترینگ نهادی + جعبه‌ابزار فنی.
\section{مقدمه}

% TODO ۲-۲ تعاریف، اصول و مبانی نظری
% این بخش الزامی نیست ولی برای این پروژه ضروری است — دو خواسته‌ی صریح کارفرما اینجا
% جای می‌گیرند:
\section{تعاریف، اصول و مبانی نظری}

% TODO ۲-۲-۱ مسئله‌ی روزنامه‌فروش (Newsvendor) و زیان نامتقارن
% عدم‌تقارن هزینه‌ی کمبود Cu در برابر مازاد Co؛ کوانتایل بهینه، رابطه (۱):
% tau*_rho = Co/(Cu+Co) در فضای نرخ عدم‌دریافت؛ چرا فرمول اولیه‌ی سند مسئله وارونه بود
% (تصحیح مستند شده به‌عنوان دستاورد، نه اشتباه پنهان‌شده).
% منبع: doc/project-definition.md بند ۲-۲؛ ⚠️ ردیف ۳۶ doc/decision_log.md
\subsection{مسئله‌ی روزنامه‌فروش و زیان نامتقارن}

% TODO ۲-۲-۲ رگرسیون کوانتایل و تابع زیان Pinball
% رابطه (۲)؛ معادل‌سازی کمینه‌ساز Pinball با کوانتایل tau؛ تفاوت کوانتایل با میانگین
% (⚠️ تناقض ظاهری log_res — ردیف ۵۱)؛ تله‌ی تجمیع «جمع کوانتایل‌ها != کوانتایل جمع».
% منبع: doc/Literature_Review.md بند ۲-۲، ۲-۳-۱
\subsection{رگرسیون کوانتایل و پیش‌بینی احتمالاتی}

% TODO ۲-۲-۳ کالیبراسیون، پوشش و بیش‌پراکندگی
% پوشش حاشیه‌ای در برابر شرطی؛ Conformal/CQR و Mondrian؛ نقض فرض استقلال افراد در
% تجمیع سطح فرد (F08) و پیامدش برای کوانتایل خوش‌بینانه.
\subsection{کالیبراسیون و بیش‌پراکندگی}

% TODO ۲-۲-۴ اعتبارسنجی و نشت اطلاعاتی در داده‌ی زمانی
% چرا K-fold تصادفی معتبر نیست؛ rolling-origin؛ انواع نشت (هدف، زمانی، Target Encoding).
% منبع: doc/progress/06-*.md؛ doc/leakage_audit.md
\subsection{اعتبارسنجی و نشت اطلاعاتی در داده‌ی زمانی}

% TODO ۲-۲-۵ ابزارهای آماری و استنباطی به‌کاررفته
% جدول: هر آزمون، چه می‌سنجد، فرض‌ها. دوازده مورد: ADF/KPSS، Ljung-Box، ARCH-LM،
% Kruskal-Wallis، Mann-Whitney/Wilcoxon، Levene، Spearman/Kendall، Diebold-Mariano،
% بوت‌استرپ بلوکی دوبعدی و n_eff، Cliff's delta، تصحیح چندگانه BH-FDR.
% فقط تعریف اینجا؛ نتیجه‌ی هر آزمون در فصل ۳/۴ با ارجاع به همین بند.
\subsection{ابزارهای آماری و استنباطی به‌کاررفته}

% TODO ۲-۳ مروری بر ادبیات موضوع
% چهار محور: نظری Newsvendor · صنایع آنالوگ no-show (هواپیمایی/هتل/رستوران) ·
% کترینگ نهادی/دانشگاهی · جعبه‌ابزار فنی (رقابت M5 Uncertainty). جدول فشرده‌ی ۱۵ منبع.
% خلأ پژوهشی: نبود پیشینه‌ی مستقیم ایرانی.
% منبع: doc/Literature_Review.md (کل)
\section{مروری بر ادبیات موضوع}

% TODO ۲-۴ نتیجه‌گیری فصل
\section{نتیجه‌گیری}
		% فصل ۲: مفاهیم اولیه و پیش‌زمینه
% !TeX root=../main.tex
\chapter{مدل‌سازی}

\section{مقدمه}
\label{sec:3-1}

فصل پیش نشان داد مسئله‌ی این پژوهش سه ویژگی دارد که هر روشی باید با آن‌ها کنار بیاید: تصمیم پیش از
مشاهده‌ی تقاضا گرفته می‌شود، هزینه‌ی کمبود و مازاد برابر نیست، و داده ترتیب زمانی دارد.
این فصل روشی را تشریح می‌کند که با همین سه محدودیت ساخته شد، و به پنج پرسش پاسخ می‌دهد:
چه کمیتی و در چه سطحی پیش‌بینی می‌شود، در لحظه‌ی تصمیم چه اطلاعاتی مجاز است، از این
اطلاعات چه ویژگی‌هایی ساخته شد، کدام خانواده‌های مدل و با چه رویه‌ای غربال شدند، و خروجی
مدل چگونه به عددی تبدیل می‌شود که آشپزخانه با آن کار می‌کند. سنجش نتیجه‌ی این روش موضوع
فصل بعد است.

ترتیب بندهای این فصل با ترتیب واقعی کار یکی نیست: روش پیشنهادی پیش از شرح داده می‌آید،
در حالی که در عمل کاوش داده بود که انتخاب روش را تعیین کرد. سهم بزرگ شوک مشترک روزانه
از نوسان نرخ، یا فروپاشی اثر شاخص کیفیت هوا پس از کنترل شهر، پیش از هر مدل‌سازی به دست
آمدند و هم فهرست ویژگی‌ها و هم معماری مدل را شکل دادند. هرجا بند \ref{sec:3-2} به شاهدی
از داده تکیه می‌کند، به بند \ref{sec:3-3-3} ارجاع رو به جلو داده شده است.

بند \ref{sec:3-2} روش پیشنهادی را از صورت‌بندی مسئله تا لایه‌ی تصمیم تشریح می‌کند. بند
\ref{sec:3-3} منابع داده، رویه‌ی پاکسازی و یکپارچه‌سازی، یافته‌های کاوش و ممیزی نشت را
می‌آورد، و بند \ref{sec:3-4} نشان می‌دهد همین روش چگونه به یک سامانه‌ی نرم‌افزاری
بازتولیدپذیر تبدیل شد.

\section{روش پیشنهادی برای پیش‌بینی نرخ عدم‌دریافت}
\label{sec:3-2}

\subsection{صورت‌بندی مسئله و سطوح داده}
\label{sec:3-2-1}

واحد تصمیم در این مسئله چهارتایی $(d,m,r,f)$ است — روز، وعده، سلف و نوع غذا — یعنی
همان سطحی که آشپزخانه برای آن یک عدد پخت می‌گیرد. متغیر هدف نرخ عدم‌دریافت
$\rho_{d,m,r,f}$ است، نسبت رزروهای دریافت‌نشده به کل رزروهای همان سلول، و نه تعداد مطلق
پرس‌های دریافت‌نشده. نرخ به سه دلیل انتخاب شد: واریانس تعداد مطلق به‌شدت به مقیاس رزرو
وابسته است و نرخ آن را در بازه‌ی $[0,1]$ نرمال می‌کند؛ الگوهای رفتاری میان سلف‌هایی که
اندازه‌شان چند برابر یکدیگر است تنها در فضای نرخ قابل‌تعمیم‌اند؛ و رابطه‌ی بدیهی «هرچه
رزرو بیشتر، عدم‌دریافت بیشتر» از دوش مدل برداشته می‌شود تا ظرفیتش صرف الگوهای پنهان‌تر
شود. آنچه برآورد می‌شود میانگین شرطی این نرخ نیست، بلکه کوانتایل $\tau$اُم توزیع شرطی
آن، $\hat\rho_{\tau,d,m,r,f}$ است، به دلیلی که بند \ref{sec:2-2-3} تشریح کرد.

همین کمیت را می‌توان در سطوح مختلفی از تجمیع مدل کرد و هر سطح ساختار متفاوتی به مدل
عرضه می‌کند: سطح سلول بیشترین جزئیات و پرنویزترین هدف را دارد، سطح عامل روز کم‌نویزترین
سیگنال را با تنها ۲۵۶ نقطه، و سطح رزرو فردی دو میلیون مشاهده با هدفی باینری. مقصود از
«شوک روز» در سطح سوم، میانگین انحراف سلف‌ها از نرخ عادی خودشان در یک روز-وعده است:
\begin{equation}
\mathrm{shock}_{d,m} \;=\; \frac{1}{|R|}\sum_{r \in R}
\big(\rho_{d,m,r} - \overline{\rho}_{m,r}\big)
\label{eq:day-shock}
\end{equation}
یعنی عددی که می‌گوید آن روز، جدا از عادت هر سلف، کل دانشگاه چقدر بیشتر یا کمتر از
همیشه غذا نگرفته است. جدول
\ref{tab:levels} پنج سطحی را که در این پژوهش ساخته شدند نشان می‌دهد. اگر سطح داده از
ابتدا برای هر مدل تثبیت شود، هر تفاوتی که بعداً میان دو مدل دیده می‌شود آمیخته‌ای از اثر
الگوریتم و اثر سطح تجمیع است و این دو از هم جدا نمی‌شوند؛ به همین دلیل سطح داده در این
پژوهش خودش یک محور آزمایش است و هر مدل روی هر سطحی که ساختارش اجازه دهد اجرا می‌شود.

\begin{table}[!ht]
\centering
\footnotesize
\singlespacing
\renewcommand{\arraystretch}{0.9}
\caption{پنج سطح داده‌ی ساخته‌شده، واحد رکورد و حجم هر کدام}
\label{tab:levels}
\vspace{2mm}
\begin{tabular}{|p{2.5cm}|p{2.1cm}|p{3.4cm}|p{4.4cm}|}
\hline
\textbf{سطح} & \textbf{واحد رکورد} & \textbf{کمیت هدف} & \textbf{حجم} \\
\hline
\lr{L1} سلول & $(d,m,r,f)$ & $\rho$ پیوسته در $[0,1]$ & $7{,}579$ رکورد؛ ۳۰ سلف، ۱۴۲ روز، ۶ شهر \\
\hline
\lr{L2} سلف$\times$وعده & $(d,m,r)$ & $\rho$ تجمیع‌شده روی نوع غذا & حدود $4{,}100$ رکورد \\
\hline
\lr{L3} عامل روز & $(d,m)$ & شوک مشترک روزانه & ۲۵۶ نقطه \\
\hline
\lr{L4} سری زمانی & سری $(m,r)$ & $\rho_t$ هر سری & ۴۱ سری؛ میانه‌ی ۱۹۷ نقطه، کمینه ۵۲ \\
\hline
\lr{L5} رزرو فردی & $(p,d,m,r,f)$ & عدم‌دریافت به‌صورت باینری & $2{,}049{,}322$ رزرو از $26{,}768$ فرد \\
\hline
\end{tabular}
\end{table}

یک سطح عمداً در این فهرست نیست. سطح $(m,r,f)$، یعنی «هر نوع غذا در هر سلف و وعده
به‌عنوان یک سری مستقل»، در نگاه اول طبیعی‌ترین گزینه به نظر می‌رسد، ولی داده $1{,}598$
سری با میانه‌ی طول تنها ۴ نقطه می‌سازد و هیچ‌کدام به ۳۰ نقطه هم نمی‌رسند؛ روی چنین
سری‌هایی نه مدل سری‌زمانی برازش می‌شود و نه ویژگی وقفه‌دار معنا دارد. پس این سطح فقط سطح
خروجی است، نه سطح آموزش. در مقابل، قاعده‌ای که کل مقایسه را ممکن می‌کند این است که هر
مدل، در هر سطحی که آموزش دیده باشد، باید خروجی‌اش را به سطح $(d,m,r,f)$ نگاشت کند و با
همان زیان \lr{Pinball} سنجیده شود؛ معیارهای سطح‌محلی مانند سطح زیر منحنی در سطح فرد یا
خطای برآورد شوک روز فقط تشخیصی‌اند و مبنای انتخاب مدل نیستند. مزیت عملی این تصمیم آن شد
که مدل منتخب مستقیماً روی همان سطحی آموزش دید که عدد در آن مصرف می‌شود، پس نه به تجزیه‌ی
سهمی نیاز افتاد و نه به هماهنگ‌سازی سلسله‌مراتبی که بند \ref{sec:2-3-3} به‌عنوان راه دوم
معرفی کرد.

\subsection{قاعده‌ی برش اطلاعاتی}
\label{sec:3-2-2}

لحظه‌ی تولید پیش‌بینی برای ناهار و شام یکی نیست (جدول \ref{tab:cutoff}) و هر چیزی که
در آن لحظه هنوز رخ نداده باشد کنار می‌رود؛ حجم رزرو روز هدف چون ۷۲ ساعت زودتر قطعی
می‌شود می‌ماند، و هر ستونی که پس از پخت مشخص می‌شود نه. پس ویژگی‌های وقفه‌دار باید آگاه
به نوع وعده ساخته شوند، و همین قید یک نامزد پرسیگنال را هم کنار گذاشت: نتیجه‌ی ناهار
روز $d$ برای پیش‌بینی شام همان روز، با احتمال نیامدن $0.159$ در برابر $0.050$ و نسبت
خطر $3.18$، که در لحظه‌ی برش شام هنوز وجود ندارد. ممیزی این مرز در بند
\ref{sec:3-3-4} و پیاده‌سازی‌اش در بند \ref{sec:3-4-2} آمده است.

\begin{table}[!ht]
\centering
\footnotesize
\singlespacing
\renewcommand{\arraystretch}{0.9}
\caption[آخرین نتیجه‌ی در دسترس در لحظه‌ی برش، به تفکیک وعده‌ی هدف]{%
آخرین نتیجه‌ی در دسترس در لحظه‌ی برش، به تفکیک وعده‌ی هدف. برای هدفِ ناهار، شام روز
$d-1$ با اینکه تاریخش پیش‌تر است هنوز سرو نشده و مجاز نیست.}
\label{tab:cutoff}
\vspace{2mm}
\begin{tabular}{|p{3.0cm}|p{3.4cm}|p{3.8cm}|}
\hline
\textbf{وعده‌ی هدف} & \textbf{لحظه‌ی برش} & \textbf{آخرین نتیجه‌ی در دسترس} \\
\hline
ناهار روز $d$ & ساعت ۱۵ روز $d-1$ & ناهار روز $d-1$ \\
\hline
شام روز $d$ & ساعت ۲۳ روز $d-1$ & شام روز $d-1$ \\
\hline
\end{tabular}
\end{table}

\subsection{مهندسی ویژگی}
\label{sec:3-2-3}

ماتریس ویژگی نهایی صد ستون دارد که در هشت دسته سازمان یافته‌اند؛ جدول \ref{tab:features}
سرفصل دسته‌ها را با مبنای زمانی هر کدام می‌آورد و رجیستری کامل در پیوست است. برای هر
ویژگی شاهدی از کاوش داده یا دانش دامنه آورده شد: با $5{,}684$ ردیف آموزش، هر ستون
اضافه سهمی از ظرفیت مدل را می‌گیرد و ستونی که فقط «شاید مفید باشد» بیش از آنکه اطلاعات
بیاورد فضای برازش را پر می‌کند.

\begin{table}[!ht]
\centering
\footnotesize
\singlespacing
\renewcommand{\arraystretch}{0.9}
\caption{هشت دسته‌ی ویژگی و مبنای زمانی هر کدام}
\label{tab:features}
\vspace{2mm}
\begin{tabular}{|p{2.5cm}|p{5.6cm}|p{4.2cm}|}
\hline
\textbf{دسته} & \textbf{نمونه} & \textbf{مبنای زمانی} \\
\hline
تقویمی & روز هفته، فاصله تا تعطیلی بعدی، طول بلوک تعطیلات، هفته‌ی نیم‌سال، دوره‌ی امتحانات & از پیش معلوم \\
\hline
هویتی & سلف، وعده، نوع غذا، شهر، خوابگاهی‌بودن سلف & ثابت \\
\hline
مقیاس رزرو & لگاریتم رزرو سلول، لگاریتم رزرو کل روز، نسبت رزرو به تاریخچه & معلوم در لحظه‌ی برش \\
\hline
نرخ تاریخی & نرخ انبساطی سلول، نرخ کوچک‌شده‌ی سلول$\times$روز هفته، نرخ تاریخی هر غذا & فقط تا لحظه‌ی برش \\
\hline
عامل روز & شوک روز با وقفه‌ی یک، دو و هفت وعده‌ی هم‌نوع؛ میانگین متحرک شوک & آگاه به نوع وعده \\
\hline
جَوّی & نوع بارش، کمینه‌ی دما، روز برفی & مقدار روز هدف \\
\hline
برهم‌کنش & روز هفته$\times$نوع غذا، شهر$\times$وعده، ترکیب رزروکنندگان$\times$شوک روز & معیار \lr{$\Delta$AIC} \\
\hline
ترکیب رزروکنندگان\footnotemark{} & میانگین و صدک ۹۰ نرخ تاریخی افرادی که آن سلول را رزرو کرده‌اند & تاریخچه تا لحظه‌ی برش \\
\hline
\end{tabular}
\end{table}
\LTRfootnotetext{Cohort features}

دسته‌ی «نرخ تاریخی» همان رمزگذاری هدف\LTRfootnote{Target encoding} است: یک متغیر
دسته‌ای مانند نام سلف، به‌جای اینکه به چند ده ستون صفر و یک باز شود، با میانگین متغیر
هدف در همان دسته جایگزین می‌شود، یعنی «این سلف تا امروز به‌طور میانگین چه نرخی داشته
است». همین کار آن را به پرخطرترین دسته هم تبدیل می‌کند، چون اگر آن میانگین روی کل
داده‌ی آموزش گرفته شود، نرخ فرداهای هر سلف در ویژگی امروزش می‌نشیند و اطلاع آینده به
گذشته نشت می‌کند. اینجا همه‌ی این ویژگی‌ها به‌صورت انبساطی و تنها تا لحظه‌ی برش هر ردیف محاسبه
شدند، و برای سلول‌های کم‌حجم که میانگینشان روی چند مشاهده بی‌ثبات است، برآورد به سمت
میانگین کل کوچک شد\LTRfootnote{Shrinkage}~\cite{EfronMorris75stein}. دسته‌ی جَوّی تنها
دسته‌ای است که مقدار روز هدف را به کار می‌برد؛ این کار نقض قاعده‌ی برش نیست چون
پیش‌بینی هواشناسی در لحظه‌ی برش موجود است، ولی در استقرار باید مقدار پیش‌بینی‌شده جای
مقدار واقعی را بگیرد و این تفاوت در ارزیابی لحاظ نشده است.

سه نامزد دیگر هم بررسی و رد شدند. شاخص کیفیت
هوا با همبستگی خام $+0.277$ نامزد قوی به نظر می‌رسید ولی این همبستگی بین‌شهری بود و
درون تهران اثری باقی نمی‌ماند ($p=0.65$)؛ قیمت غذا اثر باقی‌مانده‌ی حدود $0.001$ داشت؛ و
ترجیح غذایی شخصی، با پایداری $0.239$ برای جفت (فرد، غذا) در برابر $0.431$ برای خود فرد،
چیزی بیش از تاریخچه‌ی کلی همان فرد نمی‌گفت.

مهم‌تر از تک‌تک این‌ها، مقایسه‌ی شش مجموعه‌ی ویژگی در شکل \ref{fig:featuresets} نشان
می‌دهد افزودن ستون بیشتر لزوماً کمک نمی‌کند. معیار محور عمودی $R^2$ خارج‌نمونه است و
ارقام با یک مدل ثابت (\lr{Gradient Boosting} با پارامتر پیش‌فرض) روی یک تقسیم زمانی
$75/25$ به دست آمده‌اند؛ هدف اجرایشان ممیزی نشت بود نه انتخاب مدل، پس فقط برای مقایسه‌ی نسبی مجموعه‌ها با یکدیگر
معتبرند و مقدار مطلقشان کران پایینی محافظه‌کارانه است، چون بازه‌ی آزمون این تقسیم
غیرعادی‌ترین بخش داده است. با همین قید، دو چیز روشن می‌شود. مجموعه‌ی وقفه‌ای، با ۴۳
ویژگی که همه از تاریخچه‌ی خود سلف می‌آیند، از مجموعه‌ی صرفاً تقویمی با ۳۱ ویژگی عقب
می‌ماند — یعنی آنچه نرخ فردا را توضیح می‌دهد بیشتر شوک مشترک آن روز است تا عادت خود
سلف. و دو مجموعه‌ی بزرگ‌تر، با ۶۲ و ۷۲ ویژگی، از مجموعه‌ی ۴۷تایی عامل روز پایین‌تر
می‌مانند، که با $5{,}684$ ردیف آموزش نشانه‌ی بیش‌برازش است؛ پس انتخاب ویژگی کاری است که
باید داخل هر بخش اعتبارسنجی تکرار شود، نه یک‌بار روی کل داده.

\begin{figure}[!ht]
\centering
\includegraphics[width=0.82\textwidth]{feature_sets.png}
\caption[$R^2$ خارج‌نمونه‌ی شش مجموعه‌ی ویژگی در برابر تعداد ویژگی]{%
$R^2$ خارج‌نمونه‌ی شش مجموعه‌ی ویژگی در برابر تعداد ویژگی هر مجموعه، با یک مدل
\lr{Gradient Boosting} ثابت روی تقسیم زمانی $75/25$. ستون تیره بهترین مجموعه است. رابطه
یکنوا نیست: بزرگ‌ترین مجموعه بهترین نیست، و مجموعه‌ی وقفه‌ای با ۴۳ ویژگی از مجموعه‌ی
تقویمی با ۳۱ ویژگی پایین‌تر می‌ماند.}
\label{fig:featuresets}
\end{figure}

\subsection{خانواده‌های مدل و قیف غربالگری}
\label{sec:3-2-4}

سیزده خانواده‌ی مدل برای این مسئله تعریف شد، از رگرسیون خطی و کوانتایل تا مدل‌های
درختی، سری‌زمانی، بیزی سلسله‌مراتبی و شبکه‌ی عصبی؛ پنج مدلی که از این میان به صدر جدول
نتایج رسیدند در بند \ref{sec:2-5} معرفی شده‌اند. هر خانواده روی هر سطح داده معنا ندارد: مدل
\lr{ARIMA} روی سطح سلول که ترتیب زمانی درونش تعریف‌شده نیست برازش نمی‌شود، و مدل واریانس
شرطی روی هدف باینری سطح فرد بی‌معناست. جدول \ref{tab:applicability} این ماتریس را پیش
از شروع اجرا تثبیت می‌کند تا برای هر خانه‌ی کنارگذاشته‌شده یک دلیل فنی ثبت شود و در پایان
معلوم باشد کدام ترکیب‌ها آزموده نشده‌اند و چرا.

\newcounter{fnfamily}
\newcommand{\fnfam}{\stepcounter{fnfamily}\footnotemark[\value{fnfamily}]}
\newcommand{\fnfamtext}[1]{\stepcounter{fnfamily}\LTRfootnotetext[\value{fnfamily}]{#1}}
\setcounter{fnfamily}{\value{footnote}}

\begin{table}[!ht]
\centering
\footnotesize
\singlespacing
\renewcommand{\arraystretch}{0.9}
\caption[ماتریس کاربردپذیری خانواده‌ی مدل در برابر سطح داده]{%
ماتریس کاربردپذیری خانواده‌ی مدل در برابر سطح داده. نشان $\bullet$ یعنی خانواده روی آن
سطح اجرا می‌شود، $\circ$ یعنی تنها با شرط یا تقریب (نمونه‌گیری، کاهش بعد، یا افزودن اثر
ثابت)، و $\times$ یعنی ساختار مدل با آن سطح ناسازگار است.}
\label{tab:applicability}
\vspace{2mm}
\begin{tabular}{|p{4.4cm}|c|c|c|c|c|}
\hline
\textbf{خانواده} & \textbf{\lr{L1}} & \textbf{\lr{L2}} & \textbf{\lr{L3}} & \textbf{\lr{L4}} & \textbf{\lr{L5}} \\
\hline
خطی و تعمیم‌یافته\fnfam & $\bullet$ & $\bullet$ & $\bullet$ & $\circ$ & $\bullet$ \\
\hline
درختی و تقویت گرادیانی\fnfam & $\bullet$ & $\bullet$ & $\bullet$ & $\bullet$ & $\bullet$ \\
\hline
سری‌زمانی تک‌متغیره\fnfam & $\times$ & $\circ$ & $\bullet$ & $\bullet$ & $\times$ \\
\hline
پنلی و چندمتغیره\fnfam & $\circ$ & $\circ$ & $\bullet$ & $\circ$ & $\times$ \\
\hline
واریانس شرطی\fnfam & $\times$ & $\circ$ & $\bullet$ & $\bullet$ & $\times$ \\
\hline
کرنل و فرآیند گاوسی\fnfam & $\bullet$ & $\bullet$ & $\bullet$ & $\circ$ & $\circ$ \\
\hline
شبکه‌ی عصبی\fnfam & $\bullet$ & $\bullet$ & $\circ$ & $\bullet$ & $\bullet$ \\
\hline
سلسله‌مراتبی و بیزی\fnfam & $\bullet$ & $\bullet$ & $\bullet$ & $\circ$ & $\circ$ \\
\hline
توزیعی و شمارشی\fnfam & $\bullet$ & $\bullet$ & $\bullet$ & $\circ$ & $\circ$ \\
\hline
نمونه‌محور\fnfam & $\bullet$ & $\bullet$ & $\bullet$ & $\circ$ & $\circ$ \\
\hline
تصمیم‌محور\fnfam & $\bullet$ & $\bullet$ & $\circ$ & $\times$ & $\circ$ \\
\hline
ترکیب و آشتی\fnfam & $\bullet$ & $\bullet$ & $\bullet$ & $\bullet$ & $\bullet$ \\
\hline
کالیبراسیون\fnfam & $\bullet$ & $\bullet$ & $\bullet$ & $\bullet$ & $\bullet$ \\
\hline
\end{tabular}
\end{table}
\setcounter{fnfamily}{\value{footnote}}
\fnfamtext{Linear and generalized linear models}
\fnfamtext{Tree-based and gradient boosting}
\fnfamtext{Univariate time series}
\fnfamtext{Panel and multivariate time series}
\fnfamtext{Conditional variance (GARCH family)}
\fnfamtext{Kernel methods and Gaussian processes}
\fnfamtext{Neural networks}
\fnfamtext{Hierarchical and Bayesian models}
\fnfamtext{Distributional and count regression}
\fnfamtext{Instance-based learning}
\fnfamtext{Decision-focused learning}
\fnfamtext{Ensembling and hierarchical reconciliation}
\fnfamtext{Calibration and conformal prediction}
\setcounter{footnote}{\value{fnfamily}}

پیمایش کامل این ماتریس ممکن نیست. هر خانه‌ی جدول یک مدل نیست بلکه چند مدل است، هر مدل
فضای هایپرپارامتر خودش را دارد، و روی این‌ها هفت محور آزمایش دیگر هم سوار است (از جمله
تبدیل هدف، وزن‌دهی، و اینکه مدل سراسری باشد یا به تفکیک شهر). حاصل‌ضرب این‌ها به صدها
پیکربندی می‌رسد که هر کدام باید روی پنج بخش اعتبارسنجی اجرا شود. به‌جای آن، از یک قیف
غربالگری\LTRfootnote{Screening funnel} چهارمرحله‌ای استفاده شد که جدول
\ref{tab:funnel} خلاصه‌اش می‌کند: بودجه در هر مرحله بالا می‌رود و تعداد نامزدها پایین
می‌آید، بر پایه‌ی نصف‌سازی متوالی که بند \ref{sec:2-5} معرفی کرد.

\begin{table}[!ht]
\centering
\footnotesize
\singlespacing
\renewcommand{\arraystretch}{0.95}
\caption[چهار مرحله‌ی قیف غربالگری]{%
چهار مرحله‌ی قیف غربالگری. هر مرحله سقف زمانی خودش را دارد؛ مدلی که از آن سقف بگذرد یا
روی نمونه‌ی $25\%$ داده اجرا می‌شود یا با برچسب «محدودشده به بودجه» کنار می‌رود، تا در
جدول نهایی معلوم بماند نبودنش نتیجه‌ی ضعف بوده یا نتیجه‌ی هزینه.}
\label{tab:funnel}
\vspace{2mm}
\begin{tabular}{|p{1.9cm}|p{4.6cm}|p{2.5cm}|p{3.4cm}|}
\hline
\textbf{مرحله} & \textbf{پرسش} & \textbf{بودجه} & \textbf{معیار عبور} \\
\hline
امکان‌سنجی & آیا مدل اصلاً برازش می‌شود و چقدر طول می‌کشد؟ & یک برازش پیش‌فرض، یک بخش & بدون خطا، زیر سقف زمانی \\
\hline
غربالگری & ترتیب تقریبی مدل‌ها چیست؟ & ۲۰ آزمایش، سه بخش & هیچ مدلی حذف نمی‌شود \\
\hline
بهینه‌سازی & هر مدل تا کجا می‌تواند بهتر شود؟ & بودجه‌ی کامل، پنج بخش & آزمون زوجی در برابر بهترین مدل تا آن لحظه \\
\hline
مدل‌های منتخب & رفتار مدل بیرون از نقطه‌ی عملیاتی چگونه است؟ & همه‌ی مقادیر $\tau$، پنج بخش & ورودی بند \ref{sec:4-4-1} \\
\hline
\end{tabular}
\end{table}

حذف مدل‌ها در این قیف با رتبه‌ی خام انجام نمی‌شود، بلکه با مقایسه‌ی
زوجی\LTRfootnote{Paired comparison} در برابر بهترین مدل تا آن لحظه، یعنی سنجش اختلاف
دو مدل روی همان سلول‌های اعتبارسنجی به‌جای مقایسه‌ی میانگین کلی‌شان. معیار حذف هم
حاشیه‌ی هم‌ارزی است که بند \ref{sec:2-5} تعریف کرد، یعنی کمترین اختلافی که برای تصمیم
پخت واقعاً فرق می‌کند؛ اینجا $\delta = 0.0005$ در زیان \lr{Pinball} که معادل حدود
$0.075$ پرس در هر سلول است.

ضرورت این حاشیه از خودِ داده می‌آید. در نخستین خانواده، زیان یک مدل از بخشی به بخش
دیگرِ اعتبارسنجی با انحراف معیار حدود $0.006$ نوسان می‌کرد، یعنی $1.6$ برابر کل فاصله‌ی
بهترین تا بدترین مدل همان خانواده ($0.0037$). دلیلش این است که نرخ همه‌ی سلف‌ها در یک
روز با هم بالا و پایین می‌رود، پس هر روز بیشتر به یک مشاهده می‌مانَد تا به چند ده
مشاهده‌ی مستقل؛ با هم‌بستگی درون‌روزی $0.225$، بخشی از اعتبارسنجی که چند هزار سلول دارد
در عمل ارزش حدود ۶۸ مشاهده‌ی مستقل را پیدا می‌کند. بدون حاشیه‌ی هم‌ارزی، رتبه‌بندی
مدل‌ها بیشتر این نوسان را دنبال می‌کرد تا تفاوت واقعی الگوریتم‌ها.

بودجه‌ی جست‌وجو هم به اندازه‌ی فضای هایپرپارامتر هر مدل بسته است، نه به تعداد
پارامترهایش: اگر شمار حالت‌های ممکن از بودجه کمتر باشد همان حالت‌ها یک‌به‌یک آزموده
می‌شوند و جست‌وجوی تصادفی کنار می‌رود، وگرنه مدلی با فضای دوحالته هم ده‌ها بار همان دو
حالت را تکرار می‌کند. ترتیب پیمایش هم محورمحور است و نه خانواده‌محور: نخست اثر هر محور
آزمایش با دو مدل ارزان اندازه گرفته می‌شود، محورهایی که اثرشان از $\delta$ کمتر است
روی ساده‌ترین مقدارشان قفل می‌شوند و عدد اندازه‌گیری‌شده ثبت می‌ماند، و بودجه‌ی باقی صرف
خانواده و هایپرپارامتر می‌شود. با این قیف، بیست‌وپنج مدل به ارزیابی کامل رسیدند که
پنج‌تایشان خط پایه را با اختلاف معنادار شکستند؛ ارقام و تفسیر این نتیجه در بند
\ref{sec:4-4-1} آمده است.


\subsection{کالیبراسیون و پوشش شرطی}
\label{sec:3-2-5}

مدلی که کمترین زیان \lr{Pinball} را دارد لزوماً مدلی نیست که تعهد $\tau$ را نگه می‌دارد.
زیان، خطای متوسط را می‌سنجد و پوشش، صحت خودِ وعده را؛ همان‌طور که بند
\ref{sec:2-3-4} نشان داد این دو مستقل‌اند و باید جدا سنجیده شوند. در این داده دلیل مشخصی هم برای سنجش
جداگانه‌ی پوشش وجود دارد: حتی پس از کنترل سلف و وعده و روز هفته و تاریخ، واریانس
باقی‌مانده $7.6$ برابر مقدار دوجمله‌ای است. با چنین پراکندگی‌ای هیچ کوانتایلی که از یک
فرمول بسته‌ی توزیعی بیرون بیاید قابل‌اعتماد نیست و پوشش باید تجربی تأیید شود. به همین
دلیل کالیبراسیون در این روش یک گام جدا و اجباری برای همه‌ی مدل‌های نهایی است، نه یک
پالایش اختیاری در انتهای کار.

سه ابزار برای این گام به کار رفت. پیش‌بینی همنوای کوانتایلی~\cite{Romano19cqr} مقدار
تصحیح را از بخشی از داده‌ی آموزش برآورد می‌کند که برای همین کار کنار گذاشته شده؛ این
بخش باید برشی از انتهای پنجره‌ی آموزش باشد و نه نمونه‌ای تصادفی. بازکالیبراسیون تطبیقی
(رابطه‌ی \ref{eq:aci}) همان تصحیح را هر روز به‌روز می‌کند. و کالیبراسیون
گروهی\LTRfootnote{Mondrian conformal prediction} تصحیح را برای هر گروه جدا برآورد
می‌کند؛ لازمش این است که مدلی با پوشش بسیار بالا در سلف‌های پرحجم و بسیار پایین در
سلف‌های کم‌حجم، در میانگین کالیبره به نظر می‌رسد، در حالی که تعهد باید در هر سلف برقرار
باشد. به همین دلیل هر مدل نهایی، جدا از پوشش کلی، پوشش را به تفکیک سلف، وعده، چارک
اندازه‌ی رزرو و تهران‌بودن هم گزارش می‌کند. نتیجه‌ی اجرای این گام در بند
\ref{sec:4-4-3} آمده است.

\subsection{لایه‌ی تصمیم: از پیش‌بینی به عدد پخت}
\label{sec:3-2-6}

خروجی مدل یک نرخ است و آنچه سلف لازم دارد یک عدد پرس. تبدیل میان این دو یک ضرب و یک
سقف‌گیری است:
\begin{equation}
\widehat{CookQty}_{d,m,r,f} \;=\; \left\lceil Res_{d,m,r,f} \times
\big(1-\hat\rho_{\tau,d,m,r,f}\big) \right\rceil
\label{eq:cookqty}
\end{equation}
هیچ قید عملیاتی دیگری روی این عدد اعمال نشد: ظرفیت هیچ سلفی محدودکننده نیست و پخت هم
در بسته‌ی گسسته انجام نمی‌شود، پس نه کف و سقف دستی لازم است و نه گردکردن به مضرب. عدد
رابطه‌ی \ref{eq:cookqty} مستقیماً همان چیزی است که به آشپزخانه اعلام می‌شود.

تنها پارامتر آزاد این لایه $\tau$ است و انتخابش تصمیمی اقتصادی است، نه فنی: رابطه‌ی
\ref{eq:tau-star} نسبت هزینه‌ی کمبود به مازاد را به یک $\tau$ ترجمه می‌کند و نسبت
چهاربرابری به $\tau = 0.20$ می‌رسد. تحلیل حاشیه‌ای روی خط پایه به همین نقطه می‌رسد: گام
$0.15$ به $0.20$ به‌ازای هر پرس کمبود اضافه $4.8$ پرس مازاد صرفه‌جویی می‌کند و می‌ارزد،
ولی گام بعدی این نرخ را به $3.6$ می‌رساند و دیگر نمی‌ارزد.

آنچه باید به مدیریت سلف اعلام شود نرخ کمبود مشاهده‌شده است و نه خودِ $\tau$، چون مدل در
عمل محافظه‌کارتر از ادعای نظری‌اش رفتار می‌کند: در $\tau = 0.20$ نرخ کمبود روی بخش‌های
اعتبارسنجی حدود $9.6\%$ درمی‌آید و نه $20\%$. سه سناریوی سیاستی که بر همین اهرم ساخته
شدند در بند \ref{sec:4-5-5} آمده‌اند.

\section{داده‌های مورد نیاز}
\label{sec:3-3}

\subsection{منابع داده}
\label{sec:3-3-1}

داده‌ی خام سامانه‌ی تغذیه‌ی دانشگاه تهران در سطح رزرو فردی است: هفت فایل ماهانه با مجموع
$2{,}564{,}919$ ردیف در بازه‌ی آذر ۱۴۰۲ تا خرداد ۱۴۰۳، که هر ردیفش یک رزرو منفرد را با
شناسه‌ی دانشجو، دانشکده، رشته، مقطع، خوابگاه و وضعیت نهایی آن رزرو ثبت می‌کند. سطح
تجمیعی خروجی جداگانه‌ای از سامانه نیست و در همین پژوهش از روی همین رکوردها ساخته شد:
وضعیت نهایی هر رزرو به دو حالت دریافت و عدم‌دریافت نگاشته می‌شود و شمارش آن‌ها در هر
ترکیب روز، وعده، سلف، غذا و جنسیت $15{,}968$ رکورد می‌دهد، که پس از پاکسازی بند
\ref{sec:3-3-2} و تجمیع روی جنسیت به $7{,}579$ سلول در ۱۴۲ روز سرو، ۳۰ سلف و ۶ شهر
می‌رسد.

تقویم روزانه از دو منبع ساخته شد: تعطیلات ملی و مذهبی از یک کتابخانه‌ی استاندارد تقویم
ایران، و تقویم آموزشی دانشگاه (شروع و پایان کلاس‌ها، حذف و اضافه، میان‌ترم، امتحانات و
تعطیلات بین‌ترمی) از متن مصوب معاونت آموزشی. منبع بعدی داده‌ی جوی و شاخص کیفیت هوای
روزانه است. نسخه‌ی نخست تنها هوای تهران را داشت و همین خطا بود: شش سلف پروژه بیرون
تهران‌اند و نسبت‌دادن هوای تهران به آن‌ها یک تفاوت ثابت شهری را به شکل متغیر روزانه
درمی‌آورد. نسخه‌ی اصلاح‌شده هر شش شهر را جدا می‌گیرد، و همین اصلاح بود که کاذب‌بودن اثر
ظاهری کیفیت هوا را آشکار کرد (بند \ref{sec:3-2-3}).

منبع سوم آرشیوی از ۳۹ رویداد بافتاری است که برای همین پژوهش ساخته شد و دو نیمه دارد.
نیمه‌ی اول رویدادهای خبری‌اند که دستی گردآوری و هرکدام با نشانی منبع خبری‌اش ثبت شدند:
ماه رمضان و عید فطر، بلوک نوروز، عزای عمومی پس از انفجار کرمان و پس از سقوط بالگرد
رئیس‌جمهور، دو مرحله‌ی انتخابات مجلس، بازی‌های تیم ملی و دو دربی تهران. نیمه‌ی دوم خودکار
از همان داده‌ی جوی مشتق شد، با آستانه‌ی صریح: شاخص کیفیت هوای بیشینه‌ی دست‌کم ۱۵۵ در دو
روز پیاپی، و بارش برف دست‌کم یک سانتی‌متر. این آرشیو عمداً به ویژگی تبدیل نشد، چون هر
رویداد یکتا در ۱۴۲ روز داده تنها یک مشاهده دارد و مدلی که از آن بیاموزد چیزی جز همان یک
روز را حفظ نکرده است؛ نقشش توضیح داده‌های پرت در کاوش داده بود. تنها استثنا ماه رمضان
است که یک رژیم ۲۹ روزه‌ی بدون سرو ناهار است و نه یک رویداد نقطه‌ای.

\subsection{پاکسازی و یکپارچه‌سازی}
\label{sec:3-3-2}

هفت آزمون یکپارچگی روی داده اجرا شد. آزمون‌های حسابداری — از جمله اینکه شمار رزرو هر
سلول برابر مجموع دریافت و عدم‌دریافت آن باشد — بدون یک نقض گذشتند و درستی تجمیع را
تأیید کردند. دو اصلاح لازم شد: حذف $2{,}883$ ردیف کاملاً تکراری، و ساختن کلید یکتا از
ترکیب شناسه‌ی رزرو و فایل مبدأ و نام غذا، چون شناسه به‌تنهایی میان ماه‌های مختلف یکتا
نیست. برای تجمیع هم باید نام‌ها یکدست می‌شد: با نرمال‌سازی متن فارسی و نگاشت
بازبینی‌شده، ۷۳ نام سلف به ۳۱ و ۶۸۲ نام غذا به ۷۸ مورد کانونیک رسید.

\subsection{یافته‌های کاوش داده}
\label{sec:3-3-3}

کاوش داده پیش از هر مدل‌سازی انجام شد و هم فهرست ویژگی‌ها و هم معماری مدل را تعیین کرد.
آنچه در ادامه می‌آید تصویری از خودِ داده است: هدررفت کجا و در چه کسانی متمرکز است، چه
چیزی آن را بالا می‌برد، و چه چیزی برخلاف انتظار اثری ندارد. مقدار $p$ هر آزمون پس از
تصحیح چندگانه‌ی بند \ref{sec:2-6} گزارش شده، و هر مقایسه‌ی دوگروهی یک بار هم در سطح
سلف\,$\times$\,وعده تکرار شده تا شباهت رکوردهای یک سلف نتیجه را متورم نکند؛ هرجا این
آزمون دوم عدد متفاوتی داده، هر دو آمده است.

\textbf{هدررفت در متوسط روزها نیست، در دُم است.} نرخ عدم‌دریافت کل بازه $8.05\%$ است؛
این عدد وزنی حجم است، در برابر میانگین سادهٔ سلول‌ها که $9.59\%$ می‌دهد و به یک سلف
ده‌نفره همان وزن سلف هزارنفره را می‌دهد. ولی خودِ میانگین آنچه را برای تصمیم پخت مهم است
پنهان می‌کند: در شکل \ref{fig:target} میانه $0.080$ است و صدک نودونهم $0.400$، یعنی پنج
برابر میانه، و در سوی دیگر $4.9\%$ سلول‌ها اصلاً عدم‌دریافت ندارند.

\begin{figure}[!ht]
\centering
\includegraphics[width=0.93\textwidth]{target_distribution.png}
\caption[توزیع نرخ عدم‌دریافت در $7{,}579$ سلول داده]{%
توزیع نرخ عدم‌دریافت در $7{,}579$ سلول داده: هیستوگرام (چپ) و تابع توزیع تجمعی (راست).
میانه $0.080$ است ولی دُم راست تا $1$ ادامه دارد و $4.9\%$ رکوردها دقیقاً صفرند. هر دو
ویژگی، برآورد میانگین شرطی را برای تصمیم پخت نامناسب می‌کنند.}
\label{fig:target}
\end{figure}

\textbf{هدررفت مسئله‌ی پردیس‌های تهران است.} شکل \ref{fig:city} نرخ را به تفکیک شهر
می‌دهد: $8.7\%$ در تهران در برابر $3.1\%$ در فومن و رضوانشهر. آزمون \lr{Mann-Whitney}
این تفاوت را تأیید می‌کند ($p = 2.1 \times 10^{-286}$، و در سطح سلف\,$\times$\,وعده
$p = 6.6 \times 10^{-6}$)، و وقتی واحد مقایسه خودِ سلف باشد جدایی دو گروه تقریباً کامل
است ($\delta = 0.964$، $p = 1.8 \times 10^{-5}$). این عدد بازتاب اندازه یا نوع سلف نیست:
با افزودن تنها متغیر «در تهران بودن»، توان توضیح تفاوت میان سلف‌ها از $0.169$ به $0.620$
می‌رود، در حالی که اندازه و نوع سلف دیگر معنادار نمی‌مانند. تفاوت میان خودِ سلف‌ها هم
بزرگ‌ترین اندازه‌ی اثر کل کاوش را دارد ($\eta^2 = 0.323$، \lr{Kruskal-Wallis} با $p$
عملاً صفر). توضیح محتمل این است که در پردیس دورافتاده سلف تنها گزینه‌ی غذایی است.

\begin{figure}[!ht]
\centering
\includegraphics[width=0.80\textwidth]{city_effect.png}
\caption[نرخ عدم‌دریافت به تفکیک شهر پردیس]{%
نرخ عدم‌دریافت به تفکیک شهر پردیس (وزنی حجم). تفاوت تهران با پردیس‌های دور تا $2.8$ برابر
است و بزرگ‌ترین اثر تک‌عاملی کل داده را می‌سازد. چون هر سلف دقیقاً در یک شهر است، شهر و
هویت سلف نباید هم‌زمان وارد یک مدل شوند.}
\label{fig:city}
\end{figure}

\textbf{یک اقلیت کوچک بیشتر هدررفت را می‌سازد.} عدم‌دریافت میان دانشجویان یکنواخت پخش
نشده است: طبق شکل \ref{fig:lorenz}، ده درصد پرتکرارترین افراد $38\%$ و بیست درصد $58\%$
کل موارد را می‌سازند و ضریب جینی $0.562$ است. این تمرکز پایدار هم هست — همبستگی نرخ هر
فرد در نیمه‌ی اول بازه با نیمه‌ی دومش $+0.636$ است — یعنی این اقلیت پیش از وقوع
قابل‌شناسایی است و مداخله‌ی هدفمند روی همین گروه به‌صرفه‌تر از هر سیاست سراسری است. همین
تمرکز نشان می‌دهد رفتار افراد مستقل از هم نیست: واریانس مشاهده‌شده $14.4$ برابر چیزی است
که از تصمیم‌های مستقل انتظار می‌رود، و آزمون \lr{Levene} میان چارک‌های حجم رزرو هم همین
ناهمگنی را تأیید می‌کند ($p = 2.4 \times 10^{-147}$؛ انحراف معیار چارک کم‌رزرو $3.03$
برابر چارک پررزرو).

\begin{figure}[!ht]
\centering
\includegraphics[width=0.62\textwidth]{lorenz.png}
\caption[منحنی لورنتس تمرکز عدم‌دریافت در دانشجویان]{%
منحنی لورنتس تمرکز عدم‌دریافت در دانشجویان. فاصله‌ی منحنی از خط‌چین برابری هرچه بیشتر
باشد، تمرکز بیشتر است؛ ضریب جینی $0.562$ همین فاصله را به یک عدد خلاصه می‌کند.}
\label{fig:lorenz}
\end{figure}

\textbf{ساکن خوابگاه کمتر غذایش را رها می‌کند.} نرخ ساکنان خوابگاه $8.4\%$ است در برابر
$11.3\%$ برای غیرساکنان ($p = 5.3 \times 10^{-50}$، $\delta = -0.123$). این تفاوت مصنوع
ترکیب دانشکده‌ها نیست: شکل \ref{fig:dorm} نشان می‌دهد در ۲۲ از ۲۴ دانشکده همین ترتیب
برقرار است و آزمون زوجی \lr{Wilcoxon} روی دانشکده‌ها آن را تأیید می‌کند
($p = 8.3 \times 10^{-7}$). ساده‌ترین توضیح این است که ساکن خوابگاه هر روز در محوطه‌ی
دانشگاه است و کمتر پیش می‌آید روزِ سرو آنجا نباشد.

\begin{figure}[!ht]
\centering
\includegraphics[width=0.93\textwidth]{dorm_resident.png}
\caption[نرخ عدم‌دریافت ساکنان خوابگاه در برابر غیرساکنان]{%
نرخ عدم‌دریافت ساکنان خوابگاه در برابر غیرساکنان: توزیع نرخ شخصی دو گروه (چپ) و مقایسه‌ی
دانشکده‌به‌دانشکده (راست). در نمودار راست هر نقطه یک دانشکده است و خط‌چین خط تساوی؛
۲۲ از ۲۴ دانشکده زیر آن می‌افتند، یعنی جهت اثر داخل هر دانشکده هم برقرار است.}
\label{fig:dorm}
\end{figure}

\textbf{دانشجوی دوره‌ی شهریه‌پرداز بیشتر غذا را رها می‌کند.} نرخ به تفکیک نوع دوره‌ی
تحصیلی تفاوت روشنی دارد: نوبت دوم $14.4\%$ و شبانه $13.7\%$ در برابر روزانه $10.2\%$
($p = 2.6 \times 10^{-39}$). برنامه‌ی کلاسی این دوره‌ها با روزانه یکی است، پس ساعت حضور
توضیحش نیست. آنچه این دو دوره را جدا می‌کند شهریه است: دانشجویی که خودش هزینه‌ی تحصیل را
می‌پردازد به‌طور متوسط وضع مالی بهتری دارد و صرف‌نظر کردن از یک وعده‌ی یارانه‌ای برایش
کم‌هزینه‌تر است. اندازه‌ی اثر با این حال کوچک است ($\eta^2 = 0.0077$)، پس این عامل کنار
هویت سلف و شهر نقش درجه‌دو دارد.

\textbf{ناهار پرریسک‌تر از شام است، برخلاف فرض اولیه.} فرض ابتدایی پروژه این بود که شام
نرخ بالاتری دارد؛ داده جهت عکس را نشان داد و آزمون یک‌طرفه در جهت فرض اولیه رد شد
($\delta = -0.228$). آزمون زوجی روی همان افراد — که ترکیب جمعیت را ثابت نگه می‌دارد —
همین جهت را تأیید می‌کند ($p = 2.1 \times 10^{-57}$)، پس تفاوت به خود وعده برمی‌گردد و
نه به اینکه چه کسانی شام می‌گیرند. روز هفته هم اثر مقاومی دارد و چهارشنبه بیشترین نرخ را
می‌گیرد ($\eta^2 = 0.055$؛ $p = 3.6 \times 10^{-88}$ و در سطح سلف\,$\times$\,وعده
$p = 0.0014$). در ایام امتحانات، برخلاف انتظار، نرخ \emph{کمتر} است
($\delta = -0.144$، $p = 1.6 \times 10^{-19}$ و خوشه‌ای $p = 0.035$)، در حالی که حجم
رزرو تغییری نمی‌کند — یعنی رفتار عوض شده، نه جمعیت.

\textbf{بزرگ‌ترین اهرم تقویمی، روز پیش از تعطیلی است.} شکل \ref{fig:preholiday} انحراف
نرخ از میانگین را بر حسب فاصله تا تعطیلی بعدی می‌دهد، پس از کنترل سلف و وعده و روز هفته.
روز آخر پیش از تعطیلی $3.1$ واحد درصد بالاتر می‌رود، یعنی نزدیک چهل درصد بیشتر از روز
عادی. آزمون یک‌طرفه‌ی \lr{Mann-Whitney} این جهش را تأیید می‌کند ($p = 0.0082$) و در سطح
سلف\,$\times$\,وعده هم پابرجا می‌ماند ($p = 0.0015$)؛ در مدلی که هم‌زمان بارش برف را
کنترل می‌کند، ضریبش $+0.0134$ با $p = 0.00023$ درمی‌آید. روزهای پیش‌تر اثر کوچک‌تر و
ناهموارتری دارند. چون تقویم از پیش معلوم است، این مستقیم‌ترین اهرمی است که در دست مدیریت
سلف می‌ماند.

\begin{figure}[!ht]
\centering
\includegraphics[width=0.62\textwidth]{pre_holiday.png}
\caption[انحراف نرخ عدم‌دریافت بر حسب فاصله تا تعطیلی بعدی]{%
انحراف نرخ عدم‌دریافت از میانگین بر حسب فاصله تا تعطیلی بعدی، پس از کنترل سلف، وعده و
روز هفته. میله‌ها بازه‌ی اطمینان بوت‌استرپی‌اند. اثر روی روز آخر متمرکز است و شدتش با
طول بلوک تعطیلی رابطه دارد.}
\label{fig:preholiday}
\end{figure}

\textbf{خبر بد سلف را تکان نمی‌دهد؛ قطع سرویس تکانش می‌دهد.} بازه‌ی داده چند روز
پرالتهاب دارد: عزای عمومی پس از انفجار کرمان، پنج روز عزای عمومی پس از سقوط بالگرد
رئیس‌جمهور، چهار رویداد امنیتی، سه دوره‌ی برف سنگین و یک اعتصاب اتوبوسرانی. در آن هشت
رویدادی که سلف کاملاً سرویس داد (۲۴ روز-وعده)، نرخ تفاوت قابل‌تشخیصی با روزهای عادی
ندارد: انحراف $+0.0047$ با بازه‌ی $[-0.0087,\, +0.0173]$ و $p = 0.47$. هیچ‌کدام از این
رویدادها حتی تا سه روز پس از خودشان از $0.6$ انحراف معیار شوک روزانه فراتر نمی‌روند.
سناریوی «تعطیلی ناگهانی و رزرو سوخته» هم عملاً رخ نداد: از دوازده روزِ قطع کامل سرویس،
یازده روز اصلاً رزروی نداشتند چون سرویس پیش از باز شدن رزرو برداشته شده بود، و تنها یک
جفت روز-وعده رزرو بی‌سرویس ماند (۱۶۳ رزرو، $0.01\%$ کل داده).

آنچه اثر دارد خودِ وقفه است. شکل \ref{fig:returncurve} نرخ را در روزهای بازگشت پس از هر
شکاف سرویسِ چهارروزه یا بیشتر — بلوک نوروز، تعطیلی بین‌ترمی، رمضان — نشان می‌دهد: چهار
روز نخست بازگشت نرخ را $22\%$ نسبی بالا می‌برد ($0.0933$ در برابر $0.0797$؛ انحراف وزنی
$+0.0175$ با بازه‌ی $[+0.0126,\, +0.0228]$) و روز پنجم اثر صفر می‌شود. بدترین روز، روز
\emph{دوم} است و نه روز اول. برخلاف روزهای بد که اثرشان روی افراد بی‌ثبات چند برابر
می‌شود، روز بازگشت همه را به یک اندازه بالا می‌برد (ضریب برهم‌کنش تاریخچه$\times$بازگشت
$-0.145$ با $p = 0.11$)، پس این پدیده‌ای در سطح روز است نه در سطح فرد. این الگو کاندید
یک ویژگی دوحالته‌ی «روز بازگشت» است که در ارزیابی این پژوهش آزموده نشد.

\begin{figure}[!ht]
\centering
\includegraphics[width=0.66\textwidth]{return_curve.png}
\caption[انحراف نرخ در روزهای بازگشت پس از شکاف سرویس]{%
انحراف نرخ عدم‌دریافت در پنج روز نخستِ بازگشت پس از شکاف سرویسِ چهارروزه یا بیشتر، نسبت
به روزهای عادی (خط‌چین). چهار روز اول بالای صفرند و روز پنجم به خط عادی برمی‌گردد؛
بدترین روز، روز دوم است.}
\label{fig:returncurve}
\end{figure}

\textbf{روزهایی که بدترین به‌نظر می‌رسند، بدترین نیستند.} پنج روزی که بالاترین نرخ را
دارند در شکل \ref{fig:dailyseries} با نقطه‌ی قرمز مشخص شده‌اند، و در پنل پایینِ همان شکل
هر پنج روز نزدیک کف نمودار حجم می‌افتند: روزهایی که دانشگاه تقریباً تعطیل بوده. بدترین
روز ظاهری در کل دانشگاه ۱۳۳ رزرو داشت، در برابر حدود بیست هزار رزروِ یک روز عادی. رابطه
سیستماتیک است — هرچه حجم روزانه کمتر، نرخ بالاتر (اسپیرمن $-0.352$،
$p = 1.7 \times 10^{-5}$) — پس فهرست «بدترین روزها» اگر بر پایه‌ی نرخ ساخته شود،
تصمیم‌گیر را به روزهایی می‌برد که هدررفت مطلقشان ناچیز است.

\begin{figure}[!ht]
\centering
\includegraphics[width=0.90\textwidth]{daily_series.png}
\caption[نرخ عدم‌دریافت روزانه در برابر حجم رزرو روزانه]{%
نرخ عدم‌دریافت روزانه (بالا) در برابر حجم کل رزرو همان روز (پایین). پنج روز با بالاترین
نرخ با نقطه‌ی قرمز مشخص شده‌اند و همگی در پنل پایین نزدیک صفرند: نرخ بالا در این روزها
نتیجه‌ی کم‌شماری است، نه هدررفت بزرگ.}
\label{fig:dailyseries}
\end{figure}

\textbf{اثر آلودگی هوا وجود ندارد.} در نگاه اول رابطه‌ی شاخص کیفیت هوا با نرخ عدم‌دریافت
قوی است (اسپیرمن $+0.277$، $p = 6.2 \times 10^{-134}$)، ولی شکل \ref{fig:aqi} نشان
می‌دهد این رابطه بین‌شهری است: تهران هم‌زمان آلوده‌ترین و پرنرخ‌ترین شهر است و همین دو
واقعیت مستقل، یک همبستگی کاذب می‌سازند. پس از کنترل سلف و وعده و روز هفته، همبستگی به
$-0.013$ می‌افتد و معنادار نمی‌ماند ($p = 0.25$)، و داخل تهران — جایی که هم‌آمیختگی شهر
اصلاً وجود ندارد — چیزی باقی نمی‌ماند ($p = 0.65$). آزمون تفاوت میان دسته‌های کیفیت هوا
هم اثر آستانه‌ای نشان نمی‌دهد ($p = 0.122$). هر سیاستی که روزهای آلوده را هدف بگیرد بر
همین همبستگی کاذب بنا شده است.

\begin{figure}[!ht]
\centering
\includegraphics[width=0.90\textwidth]{aqi_spurious.png}
\caption[شاخص کیفیت هوا و نرخ عدم‌دریافت، بین شهرها و داخل تهران]{%
شاخص کیفیت هوا و نرخ عدم‌دریافت: مقایسه‌ی شش شهر (چپ) و همان رابطه تنها داخل تهران
(راست). رابطه‌ای که در نمودار چپ قوی به‌نظر می‌رسد، در نمودار راست ناپدید می‌شود — نشانه‌ی
روشن اینکه شهر مخدوش‌کننده بوده است، نه آلودگی عامل.}
\label{fig:aqi}
\end{figure}

\textbf{ناهار و شام حافظه‌ی زمانی متفاوتی دارند.} در شکل \ref{fig:acf}، نرخ شام عمدتاً از
شام دیروز خبر می‌دهد ($+0.758$، $p = 2.4 \times 10^{-22}$) و به وقفه‌ی هفت‌روزه هیچ
حساسیتی ندارد ($+0.016$، $p = 0.86$)، در حالی که ناهار ریتم دوهفتگی دارد و وقفه‌های
چهارده و بیست‌وهشت روز قوی‌ترین‌اند ($+0.393$، $p = 7.4 \times 10^{-5}$). این ریتم فقط
گردش منو نیست، چون پس از حذف اثر ثابت سلف و غذا هم باقی می‌ماند. پس یک مجموعه‌ی یکسان از
ویژگی وقفه‌دار برای هر دو وعده، برای یکی‌شان بی‌ربط می‌شود و این ویژگی‌ها به تفکیک وعده
ساخته شدند.

\begin{figure}[!ht]
\centering
\includegraphics[width=0.90\textwidth]{acf_by_meal.png}
\caption[خودهمبستگی سری روزانه‌ی نرخ، به تفکیک وعده]{%
خودهمبستگی سری روزانه‌ی نرخ عدم‌دریافت، به تفکیک وعده. میله‌های پررنگ وقفه‌های معنادارند و
خط‌چین‌ها وقفه‌های ۷، ۱۴ و ۲۸ روز را نشان می‌دهند. الگوی دو وعده تقریباً هیچ اشتراکی
ندارد.}
\label{fig:acf}
\end{figure}

\textbf{بیشترِ نوسان، شوک مشترک یک روز است.} نرخ یک سلول بیش از هر چیز به وضعیت آن روز
بستگی دارد: از واریانس نرخ سلول، $7.0\%$ نویز نمونه‌گیری است، $9.8\%$ ترکیب رزروکنندگان،
و $83.2\%$ یک شوک مشترک روزانه. این شوک سراسری است — باقیمانده‌ی جفت‌سلف‌ها در یک روز
همبستگی $+0.421$ دارد و $99.5\%$ جفت‌ها مثبت‌اند — و شکل \ref{fig:dayshock} آن را در طول
بازه نشان می‌دهد. چون شوک از روزی به روز بعد پایدار می‌ماند، با اطلاعات مجاز در لحظه‌ی
برش تا $R^2$ خارج‌نمونه‌ی $0.60$ پیش‌بینی‌پذیر است؛ برای مدل یعنی اولویت اول عامل روز است
و نه ویژگی فردی.

\begin{figure}[!ht]
\centering
\includegraphics[width=0.93\textwidth]{day_shock.png}
\caption[شوک مشترک روزانه در طول بازه‌ی داده، به تفکیک وعده]{%
شوک مشترک روزانه در طول بازه‌ی داده، به تفکیک وعده: میانگین انحراف سلف‌ها از نرخ عادی
خودشان در هر روز-وعده، طبق رابطه‌ی \ref{eq:day-shock}. دو منحنی با هم بالا و پایین
می‌روند، یعنی شوک ویژگی یک سلف نیست. نقطه‌های قرمز روزهای پیش از تعطیلی‌اند و تقریباً همه
بالای صفر می‌افتند. شکاف منحنی‌ها روزهای بدون سرو است: تعطیلی بین‌ترمی بهمن، بلوک نوروز، و
نبود ناهار در ماه رمضان. بزرگ‌ترین جهش داده، روزهای آغاز رمضان در اسفند است.}
\label{fig:dayshock}
\end{figure}

\subsection{ممیزی نشت}
\label{sec:3-3-4}

ویژگی‌ای که اطلاع پس از لحظه‌ی برش را حمل کند، در اعتبارسنجی نتیجه‌ی بسیار خوبی می‌دهد و
در استقرار واقعی شکست می‌خورد، بی‌آنکه در هیچ مرحله‌ای خطایی ظاهر شود؛ پس درستی این مرز
در سه لایه آزموده شد. لایه‌ی اول ساختاری است: تاریخچه‌ی هر فرد یک بار مستقل بازسازی و با
خروجی ماژول ویژگی‌ساز مقایسه شد ($26{,}202$ ردیف از ۳۰۰ فرد، بدون یک ناسازگاری)، و
جداگانه بررسی شد که نتیجه‌ی ناهار روز $d$ وارد تاریخچه‌ی شام همان روز نشود — تعریف
متعارف تاریخچه دقیقاً همین مرز را می‌شکند و خوش‌بینی‌اش هرچند کوچک است ($0.7201$ در برابر
$0.7179$ در سطح زیر منحنی)، نقض تعریفیِ قاعده است و نه یک تقریب. لایه‌ی دوم همبستگی هر
ویژگی با هدف را می‌سنجد؛ آستانه $0.95$ بود و بیشترین مقدار میان صد ویژگی $0.473$. لایه‌ی
سوم سقف $R^2$ خارج‌نمونه است: چون $83\%$ نوسان نرخ از شوک روزانه می‌آید و حتی مدل اثر
ثابت کامل به $0.47$ می‌رسد، سقف واقع‌بینانه‌ی این مسئله میان $0.4$ و $0.5$ است و هیچ‌کدام
از شش مجموعه‌ی ویژگی از $0.14$ فراتر نرفتند؛ مقداری بالاتر از این سقف باید نشانه‌ی نشت
خوانده شود و پیش از پذیرفته‌شدن دوباره ممیزی شود، نه مهارت مدل.

\section{معماری و پیاده‌سازی سامانه}
\label{sec:3-4}

روش بند \ref{sec:3-2} تا وقتی به کد تبدیل نشود قابل‌آزمون نیست، و درستی نتیجه هم به
طراحی روش بستگی دارد و هم به پیاده‌سازی‌اش: قاعده‌ی برشی که در یک ماژول رعایت و در ماژول
دیگر فراموش شود، عدد اعتبارسنجی را بهتر می‌کند و خطایش تنها در استقرار دیده می‌شود. این
بند پنج نقطه‌ی حساس روش را تشریح می‌کند: خط لوله، قاعده‌ی برش، بخش‌های اعتبارسنجی،
قرارداد اجرای مدل‌ها، و لایه‌ی تصمیم.

\subsection{معماری نرم‌افزاری و خط لوله}
\label{sec:3-4-1}

هر عددی که در این گزارش آمده باید با یک فرمان بازتولید شود، پس هیچ منطق مؤثری فقط داخل
نوت‌بوک نماند: نوت‌بوک‌ها روایت و بررسی چشمی‌اند و کد اصلی در بسته‌ای است که هر ماژولش
مستقل اجرا می‌شود. جدول \ref{tab:srclayers} هشت لایه‌ی این بسته و مسئولیت هرکدام را
می‌آورد. مرز لایه‌ها از روی ترتیب خط لوله کشیده شده و نه از روی نوع کد، تا هر لایه یک
ورودی مشخص از لایه‌ی پیش و یک خروجی ذخیره‌شده روی دیسک داشته باشد؛ با این ترتیب هر
مرحله جدا از بقیه اجرا و بررسی می‌شود و شکست یک مدل با شکست ساخت ویژگی اشتباه گرفته
نمی‌شود.

\begin{table}[!ht]
\centering
\footnotesize
\singlespacing
\renewcommand{\arraystretch}{0.9}
\caption{هشت لایه‌ی بسته‌ی کد و مسئولیت هر کدام}
\label{tab:srclayers}
\vspace{2mm}
\begin{tabular}{|p{2.6cm}|p{10.4cm}|}
\hline
\textbf{لایه} & \textbf{مسئولیت} \\
\hline
پیکربندی & مسیرها، ثابت‌های سراسری، و تنها منبع بذر تصادفی پروژه \\
\hline
داده & خواندن فایل خام، پاکسازی، یکدست‌سازی نام سلف و غذا، تقویم آموزشی، داده‌ی جوی، آرشیو رویدادها \\
\hline
ویژگی & قاعده‌ی برش، ساخت ماتریس ویژگی و مجموعه‌های ویژگی، رجیستری خودکار، ممیزی نشت \\
\hline
اعتبارسنجی & تولید و انجماد بخش‌های زمانی و آزمون‌های خودکار پروتکل \\
\hline
معیار و خط پایه & معیارهای عملیاتی در فضای نرخ و پرس، و هشت خط پایه‌ی مقایسه \\
\hline
مدل & رجیستری خانواده‌ها، محورهای آزمایش و شیء پیکربندی، فضای هایپرپارامتر، هارنس هر مرحله، ماژول هر خانواده \\
\hline
کاوش & توابع آزمون آماری و نمودارهای بند \ref{sec:3-3-3} \\
\hline
تصمیم و گزارش & تبدیل نرخ به عدد پخت، سناریوهای سیاستی، و تولید شکل‌های گزارش \\
\hline
\end{tabular}
\end{table}

شکل \ref{fig:pipeline} همین خط لوله را از فایل خام تا عدد پخت نشان می‌دهد. دو خروجی
میانی آن هش دارند: ماتریس ویژگی و فایل بخش‌های اعتبارسنجی. هر اجرای مدل مقدار این دو
هش را کنار نتیجه‌اش ثبت می‌کند و اجرایی که هشش با مقدار ثبت‌شده یکی نباشد از جدول مقایسه
بیرون می‌رود. بدون این کار، اجراهایی که روی پردازنده‌ی گرافیکی ابری انجام شدند هیچ
شاهدی نداشتند که روی همان داده و همان بخش‌ها اجرا شده‌اند، چون آنجا داده از فضای ابری
خوانده می‌شود و نه از مخزن کد.

\begin{figure}[!ht]
\centering
\resizebox{\textwidth}{!}{%
\begin{tikzpicture}[
  font=\small,
  bx/.style={draw, rounded corners=2pt, minimum width=2.7cm, minimum height=1.05cm,
             align=center, fill=black!4},
  art/.style={bx, fill=black!14, thick},
  gate/.style={draw, rounded corners=2pt, minimum width=2.7cm, minimum height=1.05cm,
               align=center, fill=black!30},
  hs/.style={font=\scriptsize, align=center, text=black!55},
  ar/.style={->, thick}
]
% ردیف اول: راست به چپ
\node[bx]  (d0) at (12,4.4) {\rl{داده‌ی خام رزرو}\\\rl{(۷ فایل ماهانه)}};
\node[bx]  (d1) at (8,4.4)  {\rl{پاکسازی و}\\\rl{یکپارچه‌سازی}};
\node[bx]  (d2) at (4,4.4)  {\rl{داده‌ی سلولی}\\$(d,m,r,f)$};
\node[bx]  (d3) at (0,4.4)  {\rl{ساخت ویژگی}\\\rl{با قاعده‌ی برش}};
\node[bx, minimum height=0.8cm] (ext) at (0,6.4) {\rl{تقویم، جوّ، رویدادها}};
% ردیف دوم: چپ به راست
\node[art] (f1) at (0,2.2)  {\rl{ماتریس ویژگی}};
\node[gate](f2) at (4,2.2)  {\rl{ممیزی نشت}};
\node[art] (f3) at (8,2.2)  {\rl{بخش‌های}\\\rl{اعتبارسنجی}};
\node[bx]  (f4) at (12,2.2) {\rl{هارنس آموزش}\\\rl{و ارزیابی}};
\node[hs]  at (0,1.25)  {\rl{هش داده}};
\node[hs]  at (8,1.25)  {\rl{هش بخش‌ها}};
% ردیف سوم: راست به چپ
\node[bx]  (g1) at (8,0)    {\rl{کالیبراسیون}};
\node[bx]  (g2) at (4,0)    {\rl{لایه‌ی تصمیم}};
\node[art] (g3) at (0,0)    {\rl{عدد پخت}\\\rl{هر سلول}};

\draw[ar] (d0) -- (d1);
\draw[ar] (d1) -- (d2);
\draw[ar] (d2) -- (d3);
\draw[ar] (ext) -- (d3);
\draw[ar] (d3) -- (f1);
\draw[ar] (f1) -- (f2);
\draw[ar] (f2) -- (f3);
\draw[ar] (f3) -- (f4);
\draw[ar] (f4) |- (g1);
\draw[ar] (g1) -- (g2);
\draw[ar] (g2) -- (g3);
\end{tikzpicture}}
\caption[خط لوله‌ی نرم‌افزاری از داده‌ی خام تا عدد پخت]{%
خط لوله‌ی نرم‌افزاری از داده‌ی خام تا عدد پخت. دو خروجی میانیِ برچسب‌دار، ماتریس ویژگی و
فایل بخش‌های اعتبارسنجی، ورودی اجباری هر اجرای مدل‌اند و هشِ ثبت‌شده‌شان مبنای پذیرش آن
اجرا در جدول مقایسه است. کادر تیره دروازه است: تا ممیزی نشت رد نشود، مرحله‌ی بعد اجرا
نمی‌شود.}
\label{fig:pipeline}
\end{figure}

نسخه‌بندی دو لایه دارد. گیت کد را نگه می‌دارد و هر اجرای مدل شناسه‌ی \lr{commit} و مسیر
و شماره‌ی خط تابع آموزش‌دهنده‌اش را ثبت می‌کند، پس بازتولید یک نتیجه یعنی بازگشت به همان
\lr{commit} و باز کردن همان آدرس. فایل‌های خام که برای گیت بزرگ‌اند با
\glsadd{DVC}\lr{DVC} نسخه‌بندی شده‌اند و هش \lr{SHA-256} هرکدام جداگانه هم ثبت شده تا
صحت تک‌تک فایل‌ها مستقل از ابزار قابل اثبات باشد. بذر تصادفی هم یک منبع واحد دارد که هر
نقطه‌ورود اول آن را صدا می‌زند؛ لازمش این است که مدل‌های تصادفی این پروژه هرکدام با سه
بذر اجرا شدند و اگر بذر در هر ماژول جدا تنظیم می‌شد، تفاوت میان دو مدل از تفاوت میان دو
قرعه قابل تشخیص نبود.

\subsection{پیاده‌سازی قاعده‌ی برش و ویژگی‌های آگاه به وعده}
\label{sec:3-4-2}

قاعده‌ی برش بند \ref{sec:3-2-2} دو ردیف دارد و همین دوگانگی رایج‌ترین منبع خطای نشت
است. با شماره‌گذاری سراسری وعده‌ها، یعنی دادن شماره‌ی $2t$ به ناهار روز $t$ و $2t+1$ به
شام همان روز، هر دو ردیف به یک نامساوی واحد تبدیل می‌شوند:
\begin{equation}
\mathrm{seq}_{\text{مشاهده}} \;\le\; \mathrm{seq}_{\text{هدف}} - 2
\label{eq:cutoff-seq}
\end{equation}
برای هدف ناهار، آخرین وعده‌ی مجاز ناهار روز پیش است و برای هدف شام، شام روز پیش؛ در هر
دو حالت فاصله دقیقاً دو وعده درمی‌آید. به این ترتیب همه‌ی ویژگی‌های تاریخی، چه در سطح
سلول و چه در سطح فرد، از یک تابع مشترک می‌گذرند و یک آزمون واحد کل آن‌ها را ممیزی
می‌کند.

این صورت‌بندی یک خطای واقعی را آشکار کرد. تعریف متعارف تاریخچه‌ی فرد، «همه‌ی رزروهای
پیشین همان فرد» را جمع می‌زند و برای هدفِ شام، ناهار همان روز را هم شامل می‌شود، در حالی
که در لحظه‌ی برش شام آن وعده هنوز سرو نشده است. خطا در کد ساخت ویژگی دیده نمی‌شد چون
هیچ ستون ممنوعی در آن ظاهر نمی‌شد؛ پس از نوشتن قاعده به شکل رابطه‌ی \ref{eq:cutoff-seq}،
تفاوت دو تعریف قابل اندازه‌گیری شد و اندازه‌اش در بند \ref{sec:3-3-4} آمده است.

دو ابزار دیگر همین مرز را نگه می‌دارند. رجیستری ویژگی از خودِ ماتریس ساخته می‌شود و برای
هر ستون، خانواده و وابستگی زمانی و پوشش را می‌نویسد، پس با افزودن یا حذف یک ویژگی از کد
عقب نمی‌ماند. ممیزی نشت هم یک اسکریپت اجرایی است و نه بازبینی چشمی: سه لایه‌ی آن (بند
\ref{sec:3-3-4}) پیش از ورود به مرحله‌ی اعتبارسنجی اجرا می‌شود و با شکست هرکدام، مرحله‌ی
بعد شروع نمی‌شود.

\subsection{تولید و انجماد بخش‌های اعتبارسنجی}
\label{sec:3-4-3}

فایل بخش‌های اعتبارسنجی یک بار ساخته و پس از آن منجمد می‌شود. سازنده‌اش از ۱۴۲ تاریخ
یکتای ماتریس ویژگی پنج بخش با پنجره‌ی گسترشی می‌سازد، با کف ۶۰ روز آموزش، و اگر فایل از
پیش موجود باشد آن را بازنویسی نمی‌کند. دلیل این امتناع آن است که بازتولید فایل هشش را
عوض می‌کند و هر اجرایی که پیش‌تر با هش قدیمی ثبت شده باشد دیگر با اجراهای بعدی هم‌سنج
نیست؛ پس بازتولید بی‌دقت یک فایل کوچک، جدول مقایسه‌ی هزاران اجرا را بی‌اعتبار می‌کند.

سه آزمون خودکار همراه همین ماژول اجرا می‌شوند: اینکه هیچ تاریخی هم‌زمان در آموزش و آزمون
یک بخش نباشد، اینکه آزمون هر بخش کاملاً پس از آموزش همان بخش و پس از آزمون بخش پیشین
بیاید، و اینکه مقدار یک ویژگی انبساطی با بازمحاسبه روی زیرمجموعه‌ی آموزشی همان بخش تغییر
نکند. آزمون سوم همان ادعایی را می‌سنجد که بند \ref{sec:4-3-2} بر آن تکیه می‌کند: چون هر
ویژگی تابعی از گذشته‌ی خودِ ردیف است، فاصله‌ی حذفی میان آموزش و آزمون لازم نیست.

هر اجرای مدل بخش‌هایش را از همین فایل می‌خواند و نه با صدا زدن دوباره‌ی سازنده. تفاوت
ظریف ولی مهم است: سازنده روی داده‌ای که حتی یک روز با نسخه‌ی پیشین فرق داشته باشد، بی‌صدا
مرزهای متفاوتی می‌سازد و آنگاه دو اجرا که هم‌سنج به نظر می‌رسند روی دو تقسیم متفاوت
سنجیده شده‌اند.

\subsection{قرارداد تابعی مشترک مدل‌ها و هارنس اجرا}
\label{sec:3-4-4}

همه‌ی مدل‌ها یک امضای تابعی مشترک دارند: ورودی، دو زیرمجموعه‌ی آموزش و آزمون و مقدار
$\tau$؛ خروجی، برآورد $\hat\rho_\tau$ برای هر ردیف آزمون. دو زیرمجموعه ردیف‌های خام همان
بخش‌اند و نه ماتریس ویژگیِ از پیش برش‌خورده، پس انتخاب ویژگی و پیش‌پردازش و مقیاس‌بندی
داخل خود تابع و فقط روی بخش آموزش انجام می‌شود؛ اگر ماتریس یک بار از بیرون برای همه
ساخته می‌شد، هر پیش‌پردازشی که از کل داده تغذیه شود بی‌صدا به بخش آزمون سرایت می‌کرد.
خط پایه‌ها هم همین امضا را دارند و از همان هارنس صدا زده می‌شوند، پس اختلاف یک مدل با
خط پایه اختلاف دو مسیر کد نیست.

هارنس اجرا فقط یک شیء پیکربندی می‌پذیرد که موضع اجرا روی هر نُه محور آزمایش را صریح
اعلام می‌کند، پس هیچ محوری روی پیش‌فرض نانوشته نمی‌ماند. ضرورتش از یک اجرای واقعی آمد:
نخستین اجرای خانواده‌ی خطی، بی‌آنکه جایی اعلام شود، تنها یک نقطه از فضای نُه‌محوره را
پوشش داد — یک سطح داده، یک دامنه، یک مجموعه‌ی ویژگی و یک تبدیل هدف — و در ثبت نتایج چنان
دیده می‌شد که انگار کل خانواده آزموده شده است. با اعلام صریح، گزارش پوشش می‌تواند بگوید
کدام نقاط اصلاً زده نشده‌اند. نزدن یک نقطه محدودیت مشروعی است؛ نگفتنش خواننده را به
تعمیمی می‌برد که شواهدش وجود ندارد.

هر اجرا، جدا از پیش‌بینی، معیارهای عملیاتی همان بخش و مقدار متناظر خط پایه را هم محاسبه
و ثبت می‌کند. نسخه‌ی نخست هارنس فقط زمان اجرا را ثبت می‌کرد و دو خطا در آن پنهان ماند:
یک خطای علامت که مدل را وادار می‌کرد به‌جای $\tau$ کوانتایل $1-\tau$ را حل کند، و انفجار
برون‌یابی یک اسپلاین. هر دو با افزودن محاسبه‌ی زیان \lr{Pinball} در همان حلقه آشکار
شدند، بی‌آنکه هزینه‌ای اضافه شود، چون پیش‌بینی و مقدار واقعی هر دو همان‌جا در دست‌اند.

الگوریتم \ref{alg:harness} حلقه‌ی کامل را نشان می‌دهد. سقف زمانی بند \ref{sec:3-2-4}
پیش از شروع هر آزمایش تازه بررسی می‌شود و نه در پایان حلقه، چون یک آزمایش به‌تنهایی
می‌تواند ده‌ها دقیقه طول بکشد و بررسی پس از پایان آن، سقف را بی‌اثر می‌کند.

\begin{algorithm}[!ht]
\caption{حلقه‌ی تنظیم و ارزیابی یک مدل روی بخش‌های اعتبارسنجی}
\label{alg:harness}
\begin{algorithmic}[1]
\STATE \textbf{ورودی:} تابع مدل $g$، فضای هایپرپارامتر $\Theta$، بودجه‌ی $n$ آزمایش، مقدار $\tau$
\STATE $\mathcal{F} \gets$ خواندن پنج بخش از فایل منجمد، به‌همراه هش آن
\FOR{$i = 1$ تا $n$}
\IF{زمان سپری‌شده از سقف این مدل گذشته است}
\STATE ثبت برچسب «محدودشده به بودجه» و خروج از حلقه
\ENDIF
\STATE $\theta_i \gets$ پیشنهاد نمونه‌گیر بر پایه‌ی نتیجه‌ی $\theta_1,\dots,\theta_{i-1}$
\STATE گشودن یک رکورد آزمایش با اعلام صریح موضع روی هر نُه محور
\FOR{هر یک از پنج بخش در $\mathcal{F}$}
\STATE $\hat\rho \gets$ برازش $g$ با $\theta_i$ روی بخش آموزش و پیش‌بینی روی بخش آزمون
\STATE ثبت معیارهای عملیاتی این بخش و مقدار خط پایه‌ی همان بخش
\ENDFOR
\STATE $L_i \gets$ میانگین زیان \lr{Pinball} روی پنج بخش
\ENDFOR
\STATE \textbf{خروجی:} $\theta^\star = \arg\min_i L_i$ با شاهد همگرایی و پایداری بین بخش‌ها
\end{algorithmic}
\end{algorithm}

\subsection{پیاده‌سازی لایه‌ی تصمیم}
\label{sec:3-4-5}

رابطه‌ی \ref{eq:cookqty} در یک تابع واحد پیاده شده که همه‌ی ارزیابی‌های این پژوهش، از
خط پایه‌ها تا مدل نهایی، آن را صدا می‌زنند: معیارهای عملیاتی از روی عدد پخت محاسبه
می‌شوند و نه مستقیم از $\hat\rho_\tau$. پس عددی که به آشپزخانه اعلام می‌شود و عددی که
جدول‌های فصل بعد آن را می‌سنجند از یک مسیر کد بیرون می‌آیند و اختلاف میان معیار
گزارش‌شده و رفتار واقعی سامانه ممکن نیست.

تنها محافظ این تابع، محدود کردن $\hat\rho_\tau$ به بازه‌ی $[0,1]$ پیش از ضرب است.
لازم است چون مدل‌های خطی و اسپلاینی می‌توانند بیرون این بازه پیش‌بینی کنند و مقدار
منفی، عدد پختی بزرگ‌تر از کل رزرو می‌سازد. هیچ پس‌پردازش دیگری در این لایه نیست، چون
قیدهای عملیاتی که معمولاً اینجا اعمال می‌شوند در این سلف‌ها فعال نیستند (بند
\ref{sec:3-2-6}).

مدل نهایی مستقیماً روی سطح $(d,m,r,f)$ آموزش دیده، یعنی همان سطحی که هر نوع غذا جداگانه
در آن پخته می‌شود، پس نه تجزیه‌ی سهمی لازم شد و نه هماهنگ‌سازی سلسله‌مراتبی. اگر عدد کل
یک وعده هم خواسته شود، جمع ساده‌ی عددهای پخت کافی است، ولی این جمع یک جمع حسابداری از
چند تصمیم جداگانه است و نه کوانتایل $\tau$اُم تقاضای کل: طبق بند \ref{sec:2-3-3} چنین
جمعی محافظه‌کارانه‌تر درمی‌آید، یعنی ریسک کمبود کل کمتر از $\tau$ و مازاد بیشتر از حالتی
است که تعهد در سطح کل بسته شود. این تفاوت باید همراه عدد کل گزارش شود، وگرنه تعهد
$\tau$ در سطحی خوانده می‌شود که برایش برقرار نیست.

\section{خلاصه و جمع‌بندی}
\label{sec:3-5}

این فصل روش پیشنهادی، داده‌ی پشت آن، و سامانه‌ای را که هر دو را اجرا می‌کند تشریح کرد.
کمیت هدف نرخ عدم‌دریافت در سطح $(d,m,r,f)$ است و آنچه برآورد می‌شود کوانتایل $\tau$اُم
توزیع شرطی آن است و نه میانگین. سطح داده به‌جای اینکه از پیش تثبیت شود خودش یک محور
آزمایش شد، تا اثر الگوریتم از اثر تجمیع جدا بماند، و لحظه‌ی تولید پیش‌بینی به تفکیک وعده
تعریف شد، چون شام روز پیش در لحظه‌ی تصمیم ناهار هنوز سرو نشده است.

ماتریس ویژگی نهایی صد ستون در هشت دسته دارد و پشت هر ستون شاهدی از کاوش داده یا دانش
دامنه ایستاده است. مقایسه‌ی شش مجموعه‌ی ویژگی نشان داد افزودن ستون بیشتر لزوماً کمک
نمی‌کند: مجموعه‌های بزرگ‌تر از مجموعه‌ی میانی عقب ماندند و مجموعه‌ی وقفه‌ای از مجموعه‌ی
صرفاً تقویمی. سه نامزد دیگر، شاخص کیفیت هوا و قیمت غذا و ترجیح غذایی شخصی، پس از کنترل
مخدوش‌کننده‌ها اثری نداشتند و وارد نشدند. سیزده خانواده‌ی مدل روی سطوح سازگارشان با یک
قیف چهارمرحله‌ای غربال شدند و معیار حذف، حاشیه‌ی هم‌ارزی است و نه رتبه‌ی خام، چون نوسان
زیان میان بخش‌های اعتبارسنجی از کل فاصله‌ی بهترین تا بدترین مدل یک خانواده بزرگ‌تر است.

کاوش داده تصویر روشنی از هدررفت داد. بیشترِ نوسان نرخ یک سلول شوک مشترک همان روز است
($83.2\%$)، بزرگ‌ترین اثر تک‌عاملی به هویت سلف و شهر برمی‌گردد، و مؤثرترین اهرم تقویمی
روز پیش از تعطیلی است با $3.1$ واحد درصد افزایش. چند فرض اولیه هم رد شدند: نرخ ناهار از
شام بالاتر است، در ایام امتحانات نرخ کمتر می‌شود، اثر ظاهری آلودگی هوا پس از کنترل شهر
از بین می‌رود، و خبر بد سلف را تکان نمی‌دهد در حالی که وقفه‌ی سرویس تکانش می‌دهد. ممیزی
نشت هر سه لایه را گذراند و سقف واقع‌بینانه‌ی $R^2$ خارج‌نمونه‌ی این مسئله را میان $0.4$ و
$0.5$ نشان داد.

بند پایانی نشان داد این روش چگونه به کد تبدیل شد: قاعده‌ی برش به یک نامساوی واحد،
بخش‌های اعتبارسنجی به یک فایل منجمد با هش، محورهای آزمایش به یک شیء پیکربندی اجباری، و
لایه‌ی تصمیم به همان تابعی که هر ارزیابی از آن استفاده می‌کند. فصل بعد این سامانه را
اجرا می‌کند: معیار ارزیابی و خط پایه‌ها را تعریف می‌کند، نتیجه‌ی مدل‌ها را روی بخش‌های
اعتبارسنجی و پنجره‌ی قفل‌شده گزارش می‌دهد، و اثر روش را در سطح تصمیم پخت می‌سنجد.
		% فصل ۳: مدل‌سازی
% !TeX root=../main.tex
\chapter{پیاده‌سازی}

% TODO ۴-۱ مقدمه فصل — طبق بند ۱-۷ قالب: این فصل محیط و ابزارهای اجرا، معیار ارزیابی
% نتایج، ارائه و تحلیل نتایج، و ویژگی‌ها و دستاوردهای روش پیشنهادی را می‌گوید. یک جمله
% هم درباره‌ی مرز فصل ۳ و ۴ (فصل ۳ روش و پیاده‌سازی، فصل ۴ اجرا و ارزیابی).
\section{مقدمه}

\section{محیط و ابزارهای اجرا}

% TODO ۴-۲-۱ مشخصات ماشین محلی (پردازنده، هسته، حافظه، سیستم‌عامل)؛ مسیر پردازنده‌ی
% گرافیکی ابری؛ نسخه‌ی پایتون؛ محیط مجازی و فایل قفل وابستگی (چرا pip freeze و نه بازه‌ی
% نسخه). ⚠️ مشخصات سخت‌افزاری باید از سیستم واقعی خوانده شود، نه حدس.
% منبع: ردیف ۱ doc/decision_log.md؛ requirements.lock؛ Makefile
\subsection{محیط سخت‌افزاری و نرم‌افزاری}

% TODO ۴-۲-۲ جدول (۴-۱): کتابخانه‌های اصلی، نسخه، و نقش هرکدام در پروژه (فقط موارد
% مؤثر بر نتیجه، نه کل requirements.lock).
% منبع: requirements.lock
\subsection{کتابخانه‌های اصلی}

\subsection{زیرساخت ثبت آزمایش و اجرای موازی}

ثبت آزمایش با \lr{MLflow} روی یک \lr{backend} فایلی محلی انجام شد. تا پایان فاز ۷،
۲٬۲۱۲ اجرا ثبت شده است؛ مقایسه‌ی دستی این حجم غیرقابل‌دفاع می‌شد، پس هر اجرا موظف شد
اطلاعات لازم برای بازیابی و مقایسه را خودش حمل کند. نخستین لایه‌ی این تعهد در
\textbf{نام} اجراست: قالبی ثابت که خانواده، مدل، سطح داده، متغیر هدف، مجموعه‌ی ویژگی،
نقطه‌ی $\tau$، مرحله، سید و مُهر زمانی را به‌ترتیب کنار هم می‌گذارد، چنان‌که کل مختصات
یک آزمایش از خودِ نام خوانده شود. لایه‌ی دوم فیلدهای اجباری است؛ در کنار برچسب‌های
مرحله و منبع محاسبه و خانواده، دو برچسب اختصاصی ثبت می‌شوند: \lr{\texttt{model\_type}}
(کلاس یا کتابخانه‌ی واقعی الگوریتم) و \lr{\texttt{code\_ref}} (مسیر فایل، شماره‌ی خط و
نام تابع آموزش‌دهنده، که خودکار از خود تابع استخراج می‌شود). کنار برچسب \lr{commit}
گیت، این یعنی بازتولید یک اجرا به \lr{checkout} کردن همان \lr{commit} و بازکردن دقیقاً
همان آدرس خلاصه می‌شود. پارامترها شامل هش \lr{snapshot} داده، هش \lr{fold}ها، شناسه‌ی
مجموعه‌ی ویژگی، سطح داده و سیدند، و معیارها علاوه بر زمان برازش شامل معیارهای عملیاتی
کامل‌اند (تفصیل معیارها در بند ۴-۳).

مهم‌ترین قید این زیرساخت \textbf{دروازه‌ی انصاف} است. فایل \lr{fold}ها یک بار ساخته شد،
هشش در دفترچه‌ی داده ثبت و فایل منجمد اعلام شد (رویه‌ی ساختش در بند ۳-۴-۳ آمده). هر
اجرایی با هش \lr{fold} نامنطبق از جدول مقایسه حذف می‌شود؛ بدون این قید هیچ راهی برای
اثبات اینکه یک اجرای ابری واقعاً روی همان برش‌های زمانی انجام شده وجود نداشت.

اجرای محلی اولویت اول بود، اما چند خانواده روی پردازنده زمان غیرقابل‌قبول داشتند:
شبکه‌ی عصبی روی سطح رزرو فردی با بیش از دو میلیون رکورد، فرآیند گاوسی روی ماتریسی به
ابعاد ۷٬۵۷۹ در ۷٬۵۷۹، و نمونه‌گیری بیزی سنگین. برای این‌ها مسیر پردازنده‌ی گرافیکی
رایگان طراحی شد، با قیدهایی که همان قاعده‌ی «کد در \lr{\texttt{src/}} است» را حفظ
می‌کنند: کد اصلی در نوت‌بوک نوشته نمی‌شود و یک هدف \lr{\texttt{Makefile}} کل
\lr{\texttt{src/}} را با داده‌ی لازم بسته‌بندی می‌کند؛ سلول سوم هر نوت‌بوک روی هش
\lr{snapshot} و هش \lr{fold}ها \lr{assert} می‌زند و در صورت عدم تطابق متوقف می‌شود، چون
نوت‌بوک ابری داده را از فضای ذخیره‌سازی شخصی می‌خواند نه از \lr{DVC}؛ وضعیت بهینه‌سازی
هر ده \lr{trial} ذخیره می‌شود تا اجرا پس از قطع نشست ادامه‌پذیر باشد؛ و نوت‌بوک اجراشده
\textbf{حتی در صورت شکست آزمایش} به مخزن برمی‌گردد، چون شکست هم داده است. هر چهار
نوت‌بوک برنامه‌ریزی‌شده اجرا شدند و برگشتند.

چند محور از این پروژه به‌طور طبیعی مستقل‌پذیرند: هر خانواده‌ی مدل، هر سری زمانی درون
خانواده‌ی سری‌ها، هر نقطه‌ی $\tau$ و هر سید. قاعده‌ی حاکم بر اجرای موازی یک جمله است:
\textbf{هیچ دو اجرای هم‌زمان یک فایل مشترک را نمی‌نویسند}. هر اجرا نام یکتای خودش را
می‌سازد و خروجی‌اش را زیر همان نام می‌نویسد، و مسیرهای خروجی اجراهای ابری از مسیرهای
محلی جدا نگه داشته شدند. اثر عملی این قاعده قابل اندازه‌گیری بود: مرحله‌ی غربالگری
نخستین خانواده سیصد کار مستقل بود که روی هشت پردازه‌ی موازی در حدود پنجاه‌ونه دقیقه
تمام شد، و چهار نوت‌بوک پردازنده‌ی گرافیکی هم‌زمان روی چهار حساب کاربری مستقل اجرا
شدند. استثنا هم صریح اعلام شد: خانواده‌ی ترکیب مدل‌ها و لایه‌ی کالیبراسیون
\textbf{وابسته}‌اند و فقط پس از تعیین قهرمانان اجرا می‌شوند.

\section{معیار ارزیابی}

% TODO ۴-۳-۱ pinball@tau (فضای نرخ و تعداد پرس)، coverage@tau، shortage\_rate،
% waste\_reduction\_pct، RMSE/MAE؛ چرا معیار اصلی pinball است (ارجاع به بند ۲-۲-۲).
% منبع: doc/model-evaluation-metrics.md
\subsection{مجموعه‌ی معیارها}

% TODO ۴-۳-۲ پنجره‌ی Test قفل‌شده؛ اعتبارسنجی rolling-origin (ارجاع به بند ۲-۴-۱)؛
% ارجاع به بند ۳-۴-۳ برای رویه‌ی ساخت فولدها، اینجا فقط منطق پروتکل.
% منبع: doc/progress/06-*.md
\subsection{پروتکل اعتبارسنجی}

% TODO ۴-۳-۳ خط پایه‌های B1-B6 و نقش B3 به‌عنوان رقیب واقعی.
% منبع: reports/baselines.md
\subsection{خط پایه‌ها}

% TODO ۴-۳-۴ Diebold-Mariano؛ بوت‌استرپ بلوکی دوبعدی و n\_eff؛ تصحیح چندگانه (طبق
% بند ۲-۵).
\subsection{آزمون معناداری}

\section{نتایج به‌دست‌آمده}

% TODO ۴-۴-۱ جدول (۴-۲): مقایسه با CI بوت‌استرپ بلوکی دوبعدی (فقط ۸-۱۰ ردیف شاخص،
% جدول کامل به پیوست)؛ نتایج منفی و پرونده‌های بسته‌شده گزارش شوند (معماری دومرحله‌ای،
% خوشه‌بندی شهر).
% منبع: ردیف ۴۰، ۴۴-۴۸ doc/decision_log.md
\subsection{نتایج اعتبارسنجی و انتخاب نامزدها}

% TODO ۴-۴-۲ ⚠️ مهم‌ترین صداقت گزارش: قهرمان اعتبارسنجی و قهرمان Test یکی نیستند، با
% CI و توضیح علت، نه پنهان‌کاری. تأکید بر «Test فقط یک بار باز شد».
% منبع: ردیف ۴۹ doc/decision_log.md؛ reports/phase8/
\subsection{ارزیابی نهایی روی مجموعه‌ی قفل‌شده}

% TODO ۴-۴-۳ پوشش حاشیه‌ای در برابر شرطی؛ شکاف نرخ کمبود مشاهده‌شده تا tau اسمی.
% منبع: reports/phase8/
\subsection{کالیبراسیون و پوشش}

\section{تحلیل نتایج}

% TODO ۴-۵-۱ تحلیل باقیمانده؛ تجزیه‌ی خطا؛ الگوی بدترین موارد؛ پایداری/حساسیت؛
% Overfit/Underfit.
% منبع: reports/phase8/
\subsection{تحلیل خطا و بدترین موارد}

% TODO ۴-۵-۲ اهمیت ویژگی و SHAP؛ اثرات نهایی (PDP/ALE)؛ مدل جانشین قابل‌فهم.
% منبع: reports/phase9/
\subsection{تفسیرپذیری}

% TODO ۴-۵-۳ تبیین تناقض ظاهری log\_res (کوانتایل در برابر میانگین، طبق بند ۲-۳-۲).
% منبع: ردیف ۵۱ doc/decision_log.md
\subsection{اعتبارسنجی دامنه‌ای}

% TODO ۴-۵-۴ نتیجه‌ی L1/L2 در برابر L5 با CI.
% منبع: ردیف ۵۰ doc/decision_log.md
\subsection{مقایسه‌ی سطح تجمیعی در برابر سطح فرد}

% TODO ۴-۵-۵ منحنی trade-off کمبود/هدررفت؛ سناریوهای سیاستی؛ صرفه‌جویی سالانه با CI
% (نه عدد نقطه‌ای)؛ پروتکل پایلوت.
% منبع: reports/phase10/
\subsection{تحلیل کسب‌وکاری}

\section{ویژگی‌ها و دستاوردهای روش پیشنهادی}

% TODO ۴-۶-۱ مزیت‌های عملی روش: کارکرد در سطح عملیاتی (d,m,r,f)؛ کنترل‌پذیری صریح
% نقطه‌ی tau به‌عنوان اهرم سیاستی؛ صرفه‌جویی برآوردی با CI. ⚠️ فقط ادعاهایی که در بند
% ۴-۴ شاهد عددی دارند.
% منبع: reports/phase10/؛ ردیف ۳۶ doc/decision_log.md
\subsection{مزیت‌های عملی}

\subsection{بازتولیدپذیری}

آنچه در این پروژه قفل شده چهار چیز است: نسخه‌ی داده (هش \lr{SHA-256} هر فایل خام و
پردازش‌شده)، برش‌های اعتبارسنجی (فایل منجمد \lr{fold}ها و هشش، که هر اجرای نامنطبق را
از مقایسه حذف می‌کند)، سید تصادفی (منبع واحد در ماژول پیکربندی)، و مرجع کد هر نتیجه
(\lr{commit} گیت به‌همراه آدرس دقیق فایل و خط تابع آموزش‌دهنده، هر دو در برچسب‌های همان
اجرای \lr{MLflow}).

آنچه قفل \textbf{نشده} هم باید صریح گفته شود. مخزن \lr{DVC} فقط \lr{cache} محلی دارد و
مقصد راه دور برایش تنظیم نشده، یعنی بازتولید کامل روی ماشین دیگر نیازمند انتقال دستی
داده‌ی خام است؛ مخزن گیت هم مقصد راه دور ندارد و رساندن کد به محیط ابری از مسیر بسته‌ی
فشرده انجام شد. افزون بر این، بازتولید سرتاسری فعلاً زنجیره‌ای از دستورهای ماژولی است و
یک هدف واحد که کل مسیر داده‌ی خام تا جدول نتایج را با یک فرمان اجرا کند هنوز ساخته
نشده است.

% TODO: اگر بسته‌کاری ۱۱-۴ (README + هدف واحد `make all`) اجرا شد، پاراگراف بالا
% به‌روزرسانی شود. منبع: doc/progress/11-*.md بند ۱۱.۴

% TODO ۴-۶-۳ ⚠️ حداکثر نیم صفحه (تصمیم کاربر: فقط یک اشاره‌ی کوتاه). محتوا: تقسیم کار
% «انسان صاحب تصمیم و پذیرش، عامل مجری و مستندساز»؛ قرارداد قابل‌بارگذاری AGENTS.md؛
% چهار دفتر دائمی و ضمانت مکانیکی src/report_coverage.py با کد خروجی ناصفر؛ و نتیجه‌ی
% عملی: هرجا قرارداد به قید کد تبدیل شد رعایت شد، هرجا فقط سند ماند نقض شد. جدول چهار
% اشتباه فاز ۷ و باگ سیم‌کشی فیچرست به پیوست می‌روند، نه به بدنه.
% منبع: AGENTS.md؛ doc/phase7-execution-standard.md؛ ردیف ۳۷، ۴۷ doc/decision_log.md
\subsection{اجرای پروژه با کمک عامل‌های هوش مصنوعی}

% TODO ۴-۷ خلاصه و جمع‌بندی فصل
\section{خلاصه و جمع‌بندی}
		% فصل ۴: پیاده‌سازی
% !TeX root=../main.tex
\chapter{بحث و نتیجه‌گیری}

% TODO ۵-۱ مقدمه فصل — مروری اجمالی یک‌صفحه‌ای بر مراحل انجام تحقیق.
\section{مقدمه}

\section{نتیجه‌گیری}

% TODO ۵-۲-۱ جمع‌بندی
% خلاصه‌ی یافته‌های تحقیق جاری؛ روایت «در گام نخست ... سپس ...».
\subsection{جمع‌بندی}

% TODO ۵-۲-۲ نوآوری
% نوآوری تئوری: تصحیح فرمول کوانتایل بهینه در فضای نرخ · صورت‌بندی قاعده‌ی برش
% وعده‌آگاه · رد تجربی فرض استقلال افراد و پیامدش برای کوانتایل.
% نوآوری عملی: پروتکل حاکمیتی اجرای پروژه با عامل و ضمانت مکانیکی پوشش گزارش.
\subsection{نوآوری}

% TODO ۵-۲-۳ پیشنهادها
% کار آینده: داده‌ی چندترمی، پایلوت میدانی، Newsvendor چندکالایی، استقرار (API،
% زمان‌بند، داشبورد، پایش drift، بازآموزی).
% منبع: doc/WBS-food-demand-forecasting.md فاز ۱۲؛ doc/progress/10-*.md بند ۱۰-۶
\subsection{پیشنهادها}

% TODO ۵-۲-۴ محدودیت‌ها
% ⚠️ با همان دقت نتایج نوشته شود. پوشش ۵ماهه و نبود سیگنال فصلی/ترمی؛ نبود تأیید
% میدانی؛ حساسیت به تغییر سیاست؛ عدم‌قطعیت Cu؛ نقض استقلال افراد (F08)؛ کارایی
% ضعیف‌تر برای جمعیت cold-start؛ نبود زمان ثبت رزرو و داده‌ی لغو؛ محدودیت DVC remote؛
% اتکای مسیر GPU به منابع رایگان.
\subsection{محدودیت‌ها}
		% فصل ۵: جمع‌بندی، نتیجه‌گیری و پیشنهادها
% !TeX root=../main.tex
% این فایل عنوان «فصل ۶: مراجع» را خودش نمی‌سازد — \bibliography آن را از طریق
% \thebibliography (رفع‌شده در commands.tex تا شماره‌دار و با \bibname درست باشد)
% تولید می‌کند. اینجا \chapter اضافه نکن، دو تیتر «مراجع» پشت‌سرهم می‌سازد.

% TODO ۱۵ منبع doc/Literature_Review.md + منابع فنی (پکیج holidays، Open-Meteo،
% مستندات LightGBM/CatBoost/Optuna/MLflow/SHAP) در tex/MyReferences.bib به سبک IEEE
% اضافه شوند، به‌ترتیب ظهور در متن.
{
\small
\onehalfspacing
\bibliographystyle{plain-fa}
\bibliography{./tex/MyReferences}
}
		% فصل ۶: مراجع (عنوان از \thebibliography در commands.tex می‌آید)

\appendix
% فصلهای پس از این قسمت به عنوان ضمیمه خواهند آمد.

% دستورات لازم برای تبدیل «فصل آ» به «پیوست آ» در فهرست مطالب
\addtocontents{toc}{
    \protect\renewcommand\protect\cftchappresnum{\appendixname~}%
    \protect\setlength{\cftchapnumwidth}{\mylenapp}}

% دستورات لازم برای شماره‌گذاری صفحات پیوست‌ها بشکل آ-۱ (فعلا با glossaries سازگار نیست)
% \let\Chapter\chapter
%\pretocmd{\chapter}{
%  \clearpage
%  \pagenumbering{arabic}
%  \renewcommand*{\thepage}{\rl{\thechapter-\arabic{page}}}}{}{}
%%%%%%%%%%%%%%%%%%%%%%%%%%%%%%%%%%%%%

% پیوست‌های اصلی این پروژه (کارت مدل، جدول کامل آزمون‌های آماری، فهرست تصمیمات،
% دفتر حقایق و یافته‌ها، رجیستری ویژگی) طبق doc/report-outline.md اینجا اضافه می‌شوند.
% پیوست‌های آموزشیِ خودِ قالب (آشنایی با لاتک/جدول‌ونمودار/مدیریت مراجع) حذف شدند —
% آن‌ها راهنمای استفاده از کلاس بودند، نه محتوای این گزارش.

% برگرداندن شماره‌بندی صفحات فصول
% \let\chapter\Chapter
\pagenumbering{tartibi} % اول، دوم، ...
%\baselineskip=.75cm

\begin{latin}
\baselineskip=.6cm
\latinabstract
\latinTitlePage
\end{latin}
\label{LastPage}

\end{document}
